\documentclass[11pt]{article}
\usepackage[margin=1in]{geometry}
\usepackage[T1]{fontenc}
\usepackage[utf8]{inputenc}
\usepackage{lmodern,microtype,amsmath,amssymb,amsthm,mathtools,needspace}
\usepackage[colorlinks=true,linkcolor=blue,citecolor=blue,urlcolor=blue]{hyperref}
\hypersetup{pdftitle={Growing sparse parity under a globally separated neural prior},pdfauthor={Kaarel Hanni; Dmitry Vaintrob; OpenAI Codex}}
\setlength{\emergencystretch}{2em}
\newtheorem{theorem}{Theorem}[section]
\newtheorem{lemma}[theorem]{Lemma}
\newtheorem{proposition}[theorem]{Proposition}
\newtheorem{corollary}[theorem]{Corollary}
\theoremstyle{remark}\newtheorem{remark}[theorem]{Remark}
\DeclareMathOperator{\Var}{Var}
\DeclareMathOperator{\sign}{sign}
\DeclareMathOperator{\dist}{dist}
\DeclareMathOperator{\KL}{KL}
\newcommand{\E}{\mathbb E}
\newcommand{\Pp}{\mathbb P}
\title{Growing sparse parity under a globally separated neural prior}
\author{Kaarel H\"anni \and Dmitry Vaintrob \and OpenAI Codex}
\date{October 2026}
\begin{document}
\maketitle
\begin{abstract}
We study Bayesian neural networks with Gaussian weights and
hyperbolic-tangent hidden activations, conditioned on a fixed absolute
output margin at every Boolean input and on noiseless training labels.
For each fixed subcritical width exponent $0<\gamma<1$, we obtain
sample lower bounds for learning parity on a growing sparse set of
coordinates. Let $\beta=2^{-(L-1)}$, $p=2-\beta$, and
$\theta=(1-\beta/4)/s_p>1/2$, where
$s_p=2-2\arcsin(2\sqrt{p-1}/p)/\pi$.
For $k\le d^{1-\eta}$, with fixed $\eta>0$, the eventual sample
complexity is at least
\[
 \exp\!\left\{c\max\left[k^\theta\log(ed/k),
                  k\{\log(ed/k)\}^{2-1/\theta}\right]\right\}.
\]
The proof establishes failure at specified sample sizes on this scale.
It applies to the actual posterior and to every possible mixture among
tied variational minimizers. A prediction counts as successful only
when the correct label has posterior probability greater than $2/3$.
The argument combines the paired spectrum of the limiting kernel,
uniform transport of the global margin through nuisance-coordinate
averages, and a quantitative parity obstruction on random active
faces. The natural $\exp\{c k\log(ed/k)\}$ lower bound remains open.
\end{abstract}

\section{Model and statement}
Let $\mathcal X=\{-1,1\}^d$ carry uniform probability measure.
Fix $U\subseteq[d]$ with $|U|=k$ and
$f(x)=\chi_U(x)=\prod_{j\in U}x_j$.
The learner is not given $U$. For fixed depth $L\ge2$ and hidden width
$N$, define
\[
 h^0(x)=(1,x),\qquad
 h^s(x)=\tanh(W^s h^{s-1}(x))\quad(1\le s<L),\qquad
 z_N(x)=W^L h^{L-1}(x).
\]
There are no additional biases. All weights are independent centered
Gaussians, with
\begin{equation}\label{eq:weights}
 \Var W^1_{ij}=N^{1-\gamma}/d,\qquad
 \Var W^s_{ij}=N^{-\gamma}\quad(1<s<L),\qquad
 \Var W^L_i=N^{-\gamma-1}.
\end{equation}
Fix $a>0$ and $0<\gamma<1$. Condition the prior on
$|z_N(x)|>a$ for every $x\in\mathcal X$ and then on
$f(X_i)z_N(X_i)>0$ for independent uniform training inputs
$S=(X_1,\ldots,X_n)$. Write
\begin{align*}
 q_{N,S}(x)&=\Pp\{f(x)z_N(x)>0\mid
        |z_N(y)|>a\ \forall y,\ f(X_i)z_N(X_i)>0\ \forall i\},\\
 F_N(S)&=\Pp_x\{q_{N,S}(x)\le2/3\}.
\end{align*}
The width limit is always taken first: $d,k,L,a,\gamma,S$ are fixed
while $N\to\infty$.

Put $D=d+1$ and
\begin{equation}\label{eq:exponents}
 \beta=2^{-(L-1)},\quad p=2-\beta,\quad
 s_p=2-\frac2\pi\arcsin\frac{2\sqrt{p-1}}p,\quad
 \theta=\frac{1-\beta/4}{s_p},\quad \nu=2-\frac1\theta.
\end{equation}
Then $1/2<\theta<1$ and $0<\nu<1$.
For $L=2$, $\theta=0.719367150047\ldots$ and
$\nu=0.609889261784\ldots$.
Let
\[
 t=\log(ed/k),\qquad
 G_{d,k}=\max\{k^\theta t,k t^\nu\},\qquad
 D_r=\sum_{j=0}^r\binom dj.
\]

\begin{theorem}[Failure at specified sample sizes]\label{thm:main}
Fix $L\ge2$, $a>0$, and $\eta\in(0,1)$. There exist positive
constants $c,C,A$, and sufficiently large $k_0,d_0$ such that the
following holds for $d\ge d_0$ and $k_0\le k\le d^{1-\eta}$.
Choose an odd integer $r$ within fixed constant factors of
\begin{equation}\label{eq:cutoff-choice}
 c\max\{k^\theta,k t^{\nu-1}\},
\end{equation}
with the constants chosen as in the proof, put
$\rho=r/D$, $H=\log(e/\rho)$, and set
\begin{equation}\label{eq:exact-size}
 n_r=\left\lceil A D_r rH\rho^{-1/400}\right\rceil.
\end{equation}
Then
\[
 cG_{d,k}\le\log n_r\le CG_{d,k},
\]
and, for every support $U$ of size $k$ and fixed $0<\gamma<1$,
\[
 \Pp_S\left\{\liminf_{N\to\infty}F_N(S)\ge\frac1{16}\right\}
 \ge\frac34.
\]
The conditioning events have positive probability for all sufficiently
large widths at every fixed dimension and training sample.
\end{theorem}


\begin{corollary}[Depth two]
For $L=2$, the eventual sample-complexity lower bound is
\[
 \exp\!\left\{c\max\left[
 k^{0.719367150047\ldots}\log(ed/k),\
 k\{\log(ed/k)\}^{0.609889261784\ldots}\right]\right\}.
\]
The two exponents are the exact constants defined in~\eqref{eq:exponents}.
\end{corollary}

The statement concerns the particular sizes~\eqref{eq:exact-size}.
It does not assert failure at every smaller sample size. This distinction
matters because no monotonicity of fixed-target posterior risk is used.
It nevertheless gives a lower bound for the usual eventual sample
complexity: a guarantee that holds for every $n\ge m$ cannot have
$m\le n_r$. For example, no eventual guarantee with training-sample
success probability greater than $1/4$ can ensure width-limit confidence
failure strictly below $1/16$ at every $n\ge m$ when $m\le n_r$.
The theorem does not quantify a minimum width or interchange the width
and dimension limits.

\subsection{Variational formulation and proof organization}
Let $\Psi(t)=2\arcsin(t)/\pi$ and
\[
 K(x,y)=\Psi^{\circ(L-1)}\!\left(\frac{1+x\cdot y}{d+1}\right).
\]
All $L^p$ norms and Walsh coefficients below use the uniform cube
probability measure. Let $\lambda_j$ be the degree-$j$ eigenvalue of
the normalized covariance operator $2^{-d}K$. The corresponding energy
is
\[
 \|z\|_K^2=z^{\mathsf T}K^{-1}z
      =\sum_{V\subseteq[d]}\frac{\widehat z(V)^2}{\lambda_{|V|}}.
\]
Write $Q_r=K_{>r}$ and
$\tau_r=Q_r(x,x)=\sum_{j>r}\binom dj\lambda_j$.
The width-first posterior selects compatible sign functions minimizing
$\|z\|_K^2/2$ under
\begin{equation}\label{eq:obstacle}
 |z(x)|\ge a\quad(x\in\mathcal X),\qquad
 f(X_i)z(X_i)\ge a\quad(i\le n).
\end{equation}
We prove that every minimizer has test error at least $3/8$. The final
confidence conclusion is therefore independent of the posterior weights
among tied minimizers.

The main additional ingredient beyond a norm bound is transport of the
all-input margin. A minimizing field is split into a low-degree head,
a middle band, and a sufficiently high-degree tail. Near a sampled
nuisance input, the first two parts reduce to a frozen head and a common
scalar offset. A recurrence argument and angular concentration force
the radial tail profile to return to the head's sign. A center with an
incorrect head sign consequently produces a quantitative kernel-energy
witness. These witnesses control bounded, normalized wrong-side moments
on random active faces. A scalar $L^p$ Bernstein inequality and parity
influence then force the restricted head degree to be too large.

\section{The kernel rate of the actual network}
Let $M=2^d$, and identify output functions with vectors in $\mathbb R^M$.
Define
\begin{equation}\label{eq:kernel}
 \Psi(t)=\frac2\pi\arcsin t,\qquad
 K(x,y)=\Psi^{\circ(L-1)}\!\left(\frac{1+x\cdot y}{d+1}\right).
\end{equation}
The positive definiteness of $K$ is proved in the next section.
For each Boolean function $g$, put
\begin{equation}\label{eq:energy}
 E_K(g)=\min_{g(x)z(x)\ge a\ \forall x}\frac12 z^TK^{-1}z.
\end{equation}
The minimum exists since its feasible set is closed and nonempty and
the objective is coercive. Replacing the closed inequalities by strict
ones gives the same infimum, by scaling a feasible vector by $1+\epsilon$.

\begin{proposition}[Posterior concentration on minimum-energy signs]
\label{prop:rate}
For the actual network~\eqref{eq:weights}, at every fixed
$0<\gamma<1$,
\begin{equation}\label{eq:rate}
 \lim_{N\to\infty}N^{-\gamma}\log
 \mathbb P\{g(x)z_N(x)>a\ \forall x\}=-E_K(g).
\end{equation}
Consequently, after any fixed noiseless training sample, posterior
mass tends to one on the compatible functions minimizing $E_K$.
No assertion about the relative weights of tied minimizers is needed.
\end{proposition}

\begin{proof}
For hidden layer $s$, let
$Q_N^s(x,y)=N^{-1}\sum_i h_i^s(x)h_i^s(y)$.
We first establish, for each fixed $\epsilon>0$,
\begin{equation}\label{eq:freeze}
 \mathbb P\{\|Q_N^{L-1}-K\|>\epsilon\}
       \le C_\epsilon e^{-c_\epsilon N}
\end{equation}
for all sufficiently large $N$. Any matrix norm can be used, since
$M$ is fixed. Conditional on the preceding layer, the rows of the
next hidden layer are independent bounded vectors. Hoeffding's
inequality controls all their empirical second moments. Their
preactivations have covariance $N^{1-\gamma}Q_N^{s-1}$; at the first
layer the covariance is
$N^{1-\gamma}(1+x\cdot y)/d$. Because $N^{1-\gamma}\to\infty$,
the conditional second-moment map converges to the Gaussian-sign map
\[
 Q\longmapsto \left[\Psi\!\left(
 \frac{Q(x,y)}{\sqrt{Q(x,x)Q(y,y)}}\right)\right]_{x,y}.
\]
This follows by dominated convergence for a Gaussian vector, since
each scalar variance stays bounded away from zero in a neighborhood
of the limiting covariance. The convergence is uniform in such a
neighborhood: a compactness subsequence and a continuous square root
coupling reduce the assertion to the same dominated-convergence
argument. The limiting map is continuous there. Inducting through
the fixed number of layers, with smaller fixed tolerances at earlier
layers, proves~\eqref{eq:freeze}. The first-layer limit normalizes
the diagonal $(d+1)/d$, giving the kernel in~\eqref{eq:kernel}.

Conditionally on the hidden layers, the readout is exactly Gaussian:
\begin{equation}\label{eq:readout}
 Z_N\mid Q_N^{L-1}\sim\mathcal N(0,N^{-\gamma}Q_N^{L-1}).
\end{equation}
Couple it to $Y_N=N^{-\gamma/2}K^{1/2}\xi$, where
$\xi\sim\mathcal N(0,I_M)$ is independent of the hidden layers.
On a fixed sufficiently small neighborhood of the positive-definite
matrix $K$, the square-root map is Lipschitz. Hence, for $\delta>0$,
\[
 \mathbb P\{\|Z_N-Y_N\|>\delta\}
 \le C_\epsilon e^{-c_\epsilon N}
 +\mathbb P\{\|\xi\|>\delta N^{\gamma/2}/(C\epsilon)\}.
\]
First divide logarithms by $N^\gamma$ and let $N\to\infty$;
then let $\epsilon\downarrow0$. The resulting upper limit is
$-\infty$, since $\gamma<1$. The coupled vectors are therefore
exponentially equivalent at speed $N^\gamma$.

The fixed-dimensional Gaussian vectors $Y_N$ have large-deviation
rate $I(z)=z^TK^{-1}z/2$, directly from their densities. Exponential
equivalence transfers the same bounds to $Z_N$. Applying the upper
bound to a closed margin orthant and the lower bound to its interior
gives~\eqref{eq:rate}, since their infima agree. Finally, there are
only finitely many Boolean functions for fixed $d$. Dividing each
compatible orthant probability by their sum proves the posterior
concentration claim.
\end{proof}

This is a rate statement for the network itself. It does not require
relative-error approximation of rare probabilities, nor convergence
in total variation to a Gaussian posterior.


\section{Fourier level masses of the kernel}\label{app:spectrum}
Put $D=d+1$, $\beta=2^{-(L-1)}$, and $F=\Psi^{\circ(L-1)}$.
The characters $\chi_U(x)=\prod_{j\in U}x_j$ diagonalize the
normalized covariance operator $M^{-1}K$, $M=2^d$. Let $\lambda_j$
be its eigenvalue at degree $j$. We first work on the full $D$-cube
with kernel $F(u\cdot v/D)$, whose degree-$j$ eigenvalue is denoted
by $a_j$. Constants in this and the next section may depend on $L$.
Notation $b\asymp c$ means two inequalities by positive constant
multiples, uniformly in the stated ranges.

\begin{lemma}[Power coefficients and paired eigenvalues]\label{lem:power}
Write $F(z)=\sum_{m\text{ odd}}c_mz^m$. Every odd coefficient is
positive, $\sum_m c_m=1$, and
\[
 c_m\asymp\frac{\Gamma(m-\beta)}{m!}\asymp m^{-1-\beta}
 \qquad(m\ge1\text{ odd}).
\]
The full-cube eigenvalues vanish at even degrees and strictly decrease
through odd degrees. They are positive at every odd degree. On the
original $d$-cube,
\[
 \lambda_{2k}=\lambda_{2k+1}=a_{2k+1}>0,
\]
with a final unpaired eigenvalue if $d$ is even. In particular $K$
is positive definite.
\end{lemma}
\begin{proof}
The series of $\Psi$ has positive coefficients at every odd exponent,
with
\[
 [z^{2k+1}]\Psi(z)=\frac2\pi\frac{\binom{2k}{k}}{4^k(2k+1)}.
\]
Composition preserves strict positivity and oddness. Near $z=1$ in
a slit neighborhood,
\[
 \Psi(z)=1-b(1-z)^{1/2}+O((1-z)^{3/2}),
 \qquad b=\frac{2\sqrt2}{\pi}.
\]
Induction gives
$F(z)=1-A(1-z)^\beta+O((1-z)^{3\beta})$, with $A>0$;
at each composition its leading constant changes from $A_t$ to
$b\sqrt{A_t}$, starting from $A_1=b$.
These expansions admit algebraic coefficient transfer
\cite[Chapter VI]{flajolet2009}. Here are the analytic-domain details.
On the unit circle outside $\{-1,1\}$, positivity of every odd
coefficient gives $|\Psi(z)|<1$: equality in the triangle inequality
would require $z^2=1$. Hence the iterated functions extend
analytically through every compact arc away from these endpoints.
At $1$, restrict $\arg(1-z)$ to a closed proper subinterval of
$(-\pi,\pi)$. Each square-root composition reduces this angular
range, and its positive leading coefficient keeps its image in the
slit neighborhood for the next arcsine composition. Oddness gives
the identical argument at $-1$. These local sectors and the compact
arcs provide the required two-singularity indented domain.
The contributions of $1$ and $-1$ add at odd exponents, giving
\[
 c_m\sim\frac{2A}{-\Gamma(-\beta)}m^{-1-\beta}
 \quad(m\to\infty\text{ odd}).
\]
Together with positivity of the finitely many initial coefficients
and the gamma-ratio asymptotic, this proves the coefficient comparisons.
For $L=2$ the explicit formula and Stirling's bounds suffice instead.

Let $p_m(j)$ be the probability that the simple coordinate-flip walk
on the $D$-cube reaches a specified distance-$j$ vertex in $m$ steps.
Expanding $(u\cdot v/D)^m$ shows that its degree-$j$ Fourier coefficient
is $p_m(j)$. Thus $a_j=\sum_{m\text{ odd}}c_mp_m(j)$.
For the walk recursion,
\begin{align*}
 D[p_{m+1}(j)-p_{m+1}(j+2)]
 ={}&j[p_m(j-1)-p_m(j+1)]\\
 &+(D-j-2)[p_m(j+1)-p_m(j+3)].
\end{align*}
Induction from time zero gives nonincreasing probabilities through
each parity class; terms with zero boundary coefficients are omitted.
The term $m=j$ makes the odd-eigenvalue inequalities strict and proves
positivity. Even-degree eigenvalues vanish by parity.
Restricting the distinguished coordinate to $+1$ gives
$\lambda_q=a_q+a_{q+1}$, proving the claims.
\end{proof}

For odd $r$, the exact combined multiplicity of reduced degrees
$j-1,j$ is $\binom d{j-1}+\binom dj=\binom Dj$. Consequently
\begin{equation}\label{eq:paired}
 D_r=\sum_{\substack{j\le r\\j\text{ odd}}}\binom Dj,
 \quad
 \tau_r=\sum_{\substack{j>r\\j\text{ odd}}}\binom Dj a_j,
 \quad \lambda_{r+1}=a_{r+2}.
\end{equation}
The odd cutoff is essential for these exact identities.

\begin{lemma}[Sharp level and tail estimates]\label{lem:sharp}
Uniformly for odd $1\le j\le D/4$ and odd $1\le r\le D/8$,
\begin{align}
 \binom Dj a_j&\asymp j^{-1-\beta},\label{eq:level}\\
 \tau_r&\asymp r^{-\beta},\label{eq:trace}\\
 \frac{\lambda_{r+1}}{\tau_r}
 &\asymp\frac1{r\binom D{r+2}},\label{eq:ratio}\\
 s_{2,r}:=\sum_{q>r}\binom dq\lambda_q^2
 &\le\frac{C r^{-2\beta-2}}{\binom D{r+2}}.
 \label{eq:secondmoment}
\end{align}
The same level-mass estimate holds up to $j\le\rho D$ for any fixed
$\rho<1/2$, with constants depending on $\rho$.
\end{lemma}
\begin{proof}
The gamma integral and Lemma~\ref{lem:power} give, for odd $j$,
\begin{align}
 a_j&\asymp D^{-\beta}\int_0^\infty
 u^{-1-\beta}p(u)^j(1-p(u))^{D-j}\,du,
 \quad p(u)=\frac{1-e^{-2u}}2,\label{eq:poisson}\\
 &\asymp D^{-\beta}\int_0^{1/2}p^j(1-p)^{D-j}w_\beta(p)\,dp,
 \label{eq:betaweight}
\end{align}
where $w_\beta(p)=[-\frac12\log(1-2p)]^{-1-\beta}/(1-2p)$.
Indeed the gamma integral Poissonizes the walk time, and independent
Poisson coordinate counts give independent coordinate parities with
success probability $p(u)$. Since an odd-distance vertex has zero
probability at even walk times, no extra parity factor occurs.
The weight $w_\beta$ is integrable near $1/2$ and comparable to
$p^{-1-\beta}$ on $[0,b]$ for every fixed $b<1/2$.

Fix $\rho<b<1/2$ and take $j\le\rho D$. The integral on $[0,b]$
is comparable to the complete beta integral
$B(j-\beta,D-j+1)$. For the lower comparison, a beta random variable
with these parameters has mean at most $\rho$, so Markov's inequality
puts at least a fraction $1-\rho/b$ of its mass in $[0,b]$.
The identity
\[
 \binom Dj B(j-\beta,D-j+1)
 =\frac{\Gamma(D+1)\Gamma(j-\beta)}
        {\Gamma(j+1)\Gamma(D+1-\beta)}
 \asymp D^\beta j^{-1-\beta}
\]
gives the desired mass. On $[b,1/2]$, the binomial factor decreases
with $p$ and, after multiplication by $\binom Dj$, is at most
$\exp[-D\operatorname{KL}(\rho\Vert b)]$. The remaining integrable
weight makes this contribution exponentially small. This proves
\eqref{eq:level} and its stated extended range.

Sum \eqref{eq:poisson} over the odd distances above $r$. Dropping the
parity restriction and using the binomial Markov bound gives
\[
 \tau_r\le C D^{-\beta}\int_0^\infty
 u^{-1-\beta}\min\{1,Du/(r+1)\}\,du
 \le C(r+1)^{-\beta}.
\]
For the lower bound, sum \eqref{eq:level} over the odd $j$ in
$(r,2r+2]$, using a fixed cutoff strictly below $D/2$ as needed at
the endpoint. This yields $c r^{-\beta}$; bounded $r$ are absorbed
by a change of constants. Equation~\eqref{eq:ratio} follows from
\eqref{eq:level}, \eqref{eq:trace}, and \eqref{eq:paired}.

For the second moment, use paired multiplicities to write
$s_{2,r}=\sum_{j>r,\ j\text{ odd}}\binom Dj a_j^2$.
Up to an odd cutoff $b_D$ near $D/4$, its terms are at most
$Cj^{-2\beta-2}/\binom Dj$. Consecutive odd inverse binomial
coefficients have ratio
\[
 \frac{(j+1)(j+2)}{(D-j)(D-j-1)},
\]
uniformly less than a fixed constant below one in this range.
The sum is therefore bounded by a constant times its first term,
at $j=r+2$. Beyond $b_D$, monotonicity gives the bound
$a_{b_D}\tau_{b_D}\le C D^{-2\beta-1}/\binom D{b_D}$,
which is exponentially smaller than the claimed bound uniformly
for $r\le D/8$. This proves \eqref{eq:secondmoment}.
\end{proof}

Finally, two-step binomial ratios give, uniformly for odd $r\le D/8$,
\begin{equation}\label{eq:dimensionratios}
 \binom Dr\le D_r\le C\binom Dr,
 \qquad \frac{D_r}{\binom D{r+2}}\le C(r/D)^2.
\end{equation}
All constants used here can be fixed independently of a subsequently
chosen small positive limit $r/d\to\alpha$.

\begin{corollary}[Lower bounds used by the interpolation]\label{cor:lower-spectrum}
For sufficiently large $d$,
\[
 \lambda_0\ge c/d,\qquad
 \lambda_j\ge \frac{c d^{-1-\beta}}{\binom dj}\quad(0\le j\le .49d),
 \qquad v:=\lambda_{\min}(K)\ge c(d\log d)^{-\beta}.
\]
Also $v\to0$.
\end{corollary}
\begin{proof}
The first bound follows from the level estimate at full-cube degree one.
For odd reduced degree $j$, use $\lambda_j=a_j$ and
$\binom Dj/\binom dj=D/(D-j)$. For even $j$, use
$\lambda_j=a_{j+1}$ and
$\binom D{j+1}/\binom dj=D/(j+1)$.
The extended level bound with any fixed cutoff between $.49$ and $1/2$
then gives the stated uniform lower bound, including $j=0$.

For any odd full-cube degree $j\le D$, restrict the integral
\eqref{eq:poisson} to $u\ge\tfrac12\log(2D)$. In this range
$p(u)\ge(1-1/(2D))/2$ and $1-p(u)\ge1/2$, so
$p(u)^j(1-p(u))^{D-j}\ge c2^{-D}$. Integrating $u^{-1-\beta}$ gives
$a_j\ge c2^{-D}D^{-\beta}(\log D)^{-\beta}$.
Every reduced eigenvalue equals an odd $a_j$, and the matrix eigenvalues
are $M$ times the normalized ones. Thus $v\ge c(d\log d)^{-\beta}$.
Finally, for any fixed small positive cutoff fraction, the minimum
normalized eigenvalue is at most the average in its tail, whence
$v\le M\tau_r/(M-D_r)=O(d^{-\beta})$. This proves $v\to0$.
\end{proof}

\section{A globally separated interpolation comparator}
\label{sec:comparator}
The next construction bounds the energy of every minimizer. It works
for arbitrary training labels and enforces the margin on the entire cube.

\begin{lemma}[Localized interpolation]\label{lem:comparator}
There is $d_L$ such that, for $d\ge d_L$ and $1\le n\le e^{d/25}$, with
probability at least $1-e^{-c_Ld}$
over independent uniform inputs, the following holds simultaneously for
every labeling $y_i\in\{-1,1\}$ and every $A>0$.
There is a field $h$ with
\[
 |h(x)|>A\quad\text{for every }x,\qquad
 y_i h(X_i)>A\quad\text{for every }i,
 \qquad \|h\|_K^2\le C_L A^2(d+n d^\beta).
\]
For $n=0$, the same deterministic energy bound holds using a constant field.
\end{lemma}
\begin{proof}
For $z\in\mathcal X$ define the nonnegative bump
\[
 b_z(x)=20\int_{2/5}^{9/20}(1-2p)^{\operatorname{dist}(x,z)}\,dp.
\]
It satisfies $b_z(z)=1$ and $b_z(x)\le(1/5)^{\operatorname{dist}(x,z)}$.
Its degree-$j$ Walsh coefficients, apart from their translation signs,
are
\[
 \widehat b_j=20\int_{2/5}^{9/20}(1-p)^{d-j}p^j\,dp,
 \qquad
 \binom dj\widehat b_j\le\frac{20}{d+1},\qquad
 \widehat b_j\le(3/5)^d.
\]
The middle inequality follows by extending to the complete beta integral.
By Corollary~\ref{cor:lower-spectrum},
$\lambda_j\ge c_Ld^{-1-\beta}/\binom dj$ for $j\le.49d$.
Hence these degrees contribute at most
\[
 C_Ld^{1+\beta}\sum_{j\le.49d}
       \left[\binom dj\widehat b_j\right]^2\le C_Ld^\beta
\]
to $\|b_z\|_K^2$.
For $j>.49d$, the integrand is increasing in $p$ on the integration
interval, so $\widehat b_j\le.55^{d-j}.45^j$.
Let $v=\lambda_{\min}(K)$ be the smallest matrix eigenvalue.
The same corollary gives $v\ge c_L(d\log d)^{-\beta}$ and
$\lambda_j\ge v/2^d$. Thus the remaining energy is at most
\[
 \frac{2^d}{v}\sum_{j>.49d}\binom dj .55^{2(d-j)}.45^{2j}
 \le\operatorname{poly}(d)e^{-c_*d},
\]
where
\[
 c_*=-[\log2+\mathsf h(.49)+1.02\log(.55)+.98\log(.45)]>.0062,
 \qquad \mathsf h(t)=-t\log t-(1-t)\log(1-t).
\]
The exponent decreases beyond $.49d$, so the binomial entropy bound
applies. We have proved $A_d:=\|b_z\|_K^2\le C_Ld^\beta$.

Put $J(z,w)=\langle b_z,b_w\rangle_K$.
Its normalized Fourier eigenvalues are
$\nu_j=\widehat b_j^2/\lambda_j\le .72^d/v$.
Parseval and Cauchy--Schwarz give
\[
 \mathbb E_{X,Y}|J(X,Y)|
 \le\left(\sum_j\binom dj\nu_j^2\right)^{1/2}
 \le\sqrt{\nu_{\max}A_d}.
\]
For every subset $T$ of the sampled indices, including one chosen after
the sample is seen,
\[
 \left\|\sum_{i\in T}b_{X_i}\right\|_K^2
 \le nA_d+\sum_{i\ne j}|J(X_i,X_j)|.
\]
The cross sum has expectation at most
$n(n-1)\sqrt{\nu_{\max}A_d}$.
Markov's inequality at threshold $n d^\beta$ has exponentially small
failure probability for $n\le e^{.04d}$, because
$\tfrac12\log(1/.72)>.16$.
Thus all subsets simultaneously have energy at most $C_Ln d^\beta$.

Also, except with exponentially small probability, every pair of
different sample indices has inputs at Hamming distance greater than
$.2d$. A union bound costs at most
\[
 n^2\exp[-d(\log2-\mathsf h(.2))],
 \qquad \log2-\mathsf h(.2)>.19.
\]
This event excludes repeated inputs.
At any input at most one center is within distance $.1d$; the sum of
all other bumps is at most $n(.2)^{.1d}=o(1)$.
Consequently, for every selected subset of centers, its bump sum is
$1+o(1)$ at a selected center and at most $.2+o(1)$ everywhere else.

For $T=\{i:y_i=-1\}$, take
\[
 h=2A\left(1-2\sum_{i\in T}b_{X_i}\right).
\]
The preceding estimates prove its strict global and training margins
for all sufficiently large $d$.
Finally $\|1\|_K^2=1/\lambda_0\le C_Ld$ by
Corollary~\ref{cor:lower-spectrum}; the squared triangle inequality
proves the asserted energy bound. For $n=0$, use $h=2A$.
\end{proof}

\section{One-sided squared sampling with one logarithm}
\label{sec:sampling}

The sampling estimate needed below concerns a squared one-sided loss.
It follows from a metric estimate in the proof of Lee's chaining theorem
\cite[Lemma~2.12, equation~(2.20)]{Lee2023}. We spell out the comparison:
a first-power ReLU sampling theorem alone would not give this result.
Throughout this section, $C$ denotes a universal constant that may change
from line to line, and $t_-:=\max\{-t,0\}$.

\begin{lemma}[A squared Lipschitz process]\label{lem:squared-process}
Let $a_1,\ldots,a_m\in\mathbb R^D$, and set
\[
 R=\max_{i\le m}\|a_i\|_2,
 \qquad V=\sup_{\|u\|_2\le1}
       \left(\sum_{i=1}^m\langle a_i,u\rangle^2\right)^{1/2}.
\]
For each $i$, let $\psi_i:\mathbb R\to\mathbb R$ satisfy
$\psi_i(0)=0$ and $\operatorname{Lip}(\psi_i)\le1$.
For independent uniform signs $\epsilon_i$,
\begin{equation}\label{eq:squared-process}
 \mathbb E_\epsilon\sup_{\|u\|_2\le1}
 \left|\sum_{i=1}^m\epsilon_i\psi_i(\langle a_i,u\rangle)^2\right|
 \le C R\sqrt{\log(2(m+D))}\,V.
\end{equation}
\end{lemma}

\begin{proof}
Use the genuine seminorms $N_i(u)=|\langle a_i,u\rangle|$ and the
symmetric quasimetric
\[
 d_A(u,v)^2
 =\sum_{i=1}^m N_i(u-v)^2\bigl(N_i(u)^2+N_i(v)^2\bigr).
\]
The stronger intermediate estimate established in the proof of
\cite[Lemma~2.12]{Lee2023} is
\begin{equation}\label{eq:lee-metric}
 \gamma_2(T,d_A)
 \le C\bigl(\lambda\sqrt{\log(2D)}+\kappa\bigr)
       \sup_{u\in T}\left(\sum_i N_i(u)^2\right)^{1/2},
 \qquad T\subseteq B_2^D,
\end{equation}
where
\[
 \lambda=\max_i\bigl(\mathbb E N_i(g)^2\bigr)^{1/2},
 \qquad \kappa=\mathbb E\max_i N_i(g),
 \qquad g\sim\mathcal N(0,I_D).
\]
Here $\gamma_2$ is the generic-chaining functional. This is equation
(2.20) of that proof, with the energy normalization (2.17) undone.
Seminorms follow from the norm statement by replacing
$N_i(u)$ with $(N_i(u)^2+\varepsilon^2\|u\|_2^2)^{1/2}$ and then
letting $\varepsilon\downarrow0$; the resulting metrics dominate
$d_A$, and the parameters on the right converge.
For our rank-one seminorms,
$\lambda=R$ and $\kappa\le R\sqrt{2\log(2m)}$.

Let $X_u=\sum_i\epsilon_i\psi_i(\langle a_i,u\rangle)^2$.
Its subgaussian increment metric is bounded by $\sqrt2\,d_A$, because
\begin{align*}
 &\sum_i\left(\psi_i(\langle a_i,u\rangle)^2
             -\psi_i(\langle a_i,v\rangle)^2\right)^2\\
 &\quad\le\sum_i |\langle a_i,u-v\rangle|^2
       \bigl(|\langle a_i,u\rangle|+|\langle a_i,v\rangle|\bigr)^2
 \le 2d_A(u,v)^2.
\end{align*}
We used both Lipschitz continuity and
$|\psi_i(t)|\le|t|$. The subgaussian chaining inequality, together
with $X_0=0$, bounds $\mathbb E\sup X_u$ by
$C\gamma_2(B_2^D,d_A)$. Apply the same argument to $-X_u$ and add.
Equation~\eqref{eq:lee-metric} now proves
\eqref{eq:squared-process}.
\end{proof}

\begin{remark}[The normalization in the metric estimate]
The proof of \cite[Lemma~2.12]{Lee2023} passes through a coarse additive
constant in its equation (2.26), after normalizing
$\sup_T\sum_iN_i^2=1$. The homogeneous version
\eqref{eq:lee-metric} follows by retaining the scale in the two geometric
terms of that argument. Indeed, in the rank-one case $T\subseteq B_2^D$
implies $N_i(u)\le\lambda$ on $T$; writing
$\|u\|_{\mathcal N}=\max_iN_i(u)$ gives
$d_A(u,v)\le\sqrt2\|u-v\|_{\mathcal N}$ on $T$.
Thus $\operatorname{diam}_{d_A}(T)\le C\lambda$ and the volumetric
packing constant in its Corollary~2.11 can be taken to be
$C\lambda$, since $T\subseteq\lambda B_{\mathcal N}$.
The additive term is therefore $C\lambda$, absorbed by
$C(\lambda\sqrt{\log(2D)}+\kappa)$.
The growth-functional portion of the proof is unchanged.
If the packing constant is less than one, first rescale the quasimetric
to apply that corollary in its stated normalization, and then scale back.
Most importantly, throughout that argument the $N_i$ are symmetric
seminorms. No assertion that a positive part is a seminorm is used.
\end{remark}

\begin{lemma}[Uniform one-sided squared sampling]
\label{lem:squared-sampling}
Let $(Z,Y)$ be a random pair with $Z\in\mathbb R^D$, $Y\in\{-1,1\}$,
$\mathbb E ZZ^{\mathsf T}=I_D$, and $\|Z\|_2^2=D$ almost surely.
For independent copies $(Z_i,Y_i)$, let
\[
 \Delta_n=\sup_{\|u\|_2\le1}
 \left|\frac1n\sum_{i=1}^n(Y_i\langle Z_i,u\rangle)_-^2
          -\mathbb E(Y\langle Z,u\rangle)_-^2\right|,
 \qquad h_n=\sqrt{\frac{D\log(2(n+D))}{n}}.
\]
Then $\mathbb E\Delta_n\le C h_n(1+C h_n)$.
In particular, for any fixed $\varepsilon,\delta>0$ there is a finite
$A_{\varepsilon,\delta}$ such that
\[
 n\ge A_{\varepsilon,\delta}D\log(2D)
 \quad\Longrightarrow\quad
 \mathbb P\{\Delta_n>\varepsilon\}\le\delta.
\]
The labels may depend arbitrarily on the features.
\end{lemma}

\begin{proof}
Put $Q_n=n^{-1}\sum_iZ_iZ_i^{\mathsf T}$.
Conditional on the observations, apply
Lemma~\ref{lem:squared-process} to
$a_i=Y_iZ_i/\sqrt n$ and $\psi_i(t)=t_-$.
Here $R=\sqrt{D/n}$ and $V=\sqrt{\|Q_n\|_{\mathrm{op}}}$.
Symmetrization gives
\begin{equation}\label{eq:sampling-symmetrization}
 \mathbb E\Delta_n
 \le C h_n\,\mathbb E\sqrt{\|Q_n\|_{\mathrm{op}}}.
\end{equation}
We also apply the same process lemma with $\psi_i(t)=t$.
Since $\mathbb E Q_n=I_D$, symmetrization and Jensen's inequality imply,
with $A=\mathbb E\|Q_n\|_{\mathrm{op}}$,
\[
 A\le1+\mathbb E\|Q_n-I_D\|_{\mathrm{op}}
 \le1+C h_n\mathbb E\sqrt{\|Q_n\|_{\mathrm{op}}}
 \le1+C h_n\sqrt A.
\]
Solving this quadratic inequality yields $\sqrt A\le1+C h_n$,
after increasing $C$. Substitute into
\eqref{eq:sampling-symmetrization}.
Finally, at $n\ge A_{\varepsilon,\delta}D\log(2D)$ the quantity
$h_n$ can be made uniformly small by increasing
$A_{\varepsilon,\delta}$. Markov's inequality proves the probability
statement. It is a statement for each specified sample size $n$.
\end{proof}


\section{Random tail covariance matrices}\label{app:gram}
Let $Q$ be the covariance kernel retaining only degrees greater than an
odd cutoff $r$, and write $\tau=\tau_r$, $\ell=\lambda_{r+1}$, and
$s_2=s_{2,r}$. This section proves the sample event needed later, without
assuming that a Gaussian approximation remains valid after conditioning.
The entire argument concerns the exact limiting Gaussian covariance.

\begin{lemma}[Far-pair decay]\label{lem:far}
There are constants $c_0,C>0$ such that, when $r\le D/8$ is odd,
\[
 \max_{d/4\le\dist(x,y)\le3d/4}|Q(x,y)|
 \le\rho_r:=C\sqrt D\,e^{-r/8}+Ce^{-c_0D}.
\]
\end{lemma}
\begin{proof}
The full-cube Krawtchouk polynomial satisfies
\[
 \sum_{j=0}^D\mathsf K_j(h;D)z^j=(1+z)^{D-h}(1-z)^h.
\]
If $\delta D\le h\le(1-\delta)D$, put $b=\min(h,D-h)$.
On $|z|=t$ the generating function is bounded in absolute value by
\[
 (1+t^2)^b(1+t)^{D-2b}
 \le(1+t)^D\exp\!\left[-\frac{2\delta Dt}{(1+t)^2}\right].
\]
Cauchy's coefficient bound at $t=j/(D-j)$ and the uniform Stirling
lower bound for $\binom Dj$ therefore give, for $1\le j<D$,
\begin{equation}\label{eq:kraw}
 \frac{|\mathsf K_j(h;D)|}{\binom Dj}
 \le C\sqrt D\exp[-2\delta j(1-j/D)].
\end{equation}
For reduced-cube points at distance in $[d/4,3d/4]$, their augmented
full-cube distance lies in $[D/8,7D/8]$ for large $d$.
Because $r$ is odd, the exact tail is
$Q(x,y)=\sum_{j>r,\ j\text{ odd}}a_j\mathsf K_j(h;D)$.
For $r<j\le D-r$, \eqref{eq:kraw} gives a bound
$C\sqrt D e^{-r/8}$ after summation, since the total spectral mass
is one. For $j>D-r$, monotonicity and an odd cutoff $b_D$ near $D/4$
give
\[
 \sum_{j>D-r}\binom Dj a_j
 \le a_{b_D}\sum_{k<r}\binom Dk
 \le C D^{-\beta-1/2}
       \exp[-D(\mathsf h(1/4)-\mathsf h(1/8))+O(1)].
\]
Here $\mathsf h(t)=-t\log t-(1-t)\log(1-t)$ denotes the binary entropy in nats.
The positive entropy gap absorbs the polynomial factor and proves the
claim. Adjust constants for bounded $d$.
\end{proof}

\begin{lemma}[A tail Gram deviation bound]\label{lem:gram-general}
Let $X_1,\ldots,X_n$ be independent uniform cube points, with
$1\ll n\le M/4$. Let
\[
 B_{\mathrm{near}}=\#\{y:\dist(x,y)\notin[d/4,3d/4]\}
 \le 2e^{d\mathsf h(1/4)},\qquad q=2\log M.
\]
The count is independent of $x$. For every $z>0$,
\begin{align}
 &\mathbb P\{\|Q_{SS}-\tau I_n\|_{\mathrm{op}}>z\tau\}
 \notag\\
 &\quad\le\frac C{z^2}
 \left[\frac{q\rho_r}{\tau}
       +\frac{\sqrt{nq s_2}}{\tau}
       +\frac{n\ell}{\tau}
       +\frac{n(1+B_{\mathrm{near}})}M\right]^2
       +\frac{n^2(1+B_{\mathrm{near}})}M.
 \label{eq:gram-bound}
\end{align}
Here the sampled matrix includes repeated training locations.
\end{lemma}
\begin{proof}
Let $J=Q-\tau I_M$, a zero-diagonal Hermitian matrix. Form a deterministic
mask $J_{\mathrm{far}}$ retaining only pairs at distances in
$[d/4,3d/4]$. Since $Q$ is positive semidefinite with constant diagonal
$\tau$, every entry has absolute value at most $\tau$. Thus
\begin{align*}
 \|J_{\mathrm{far}}\|_{\max}&\le\rho_r,\\
 \|J_{\mathrm{far}}\|_{1\to2}^2&\le M s_2,\\
 \|J_{\mathrm{far}}\|_{\mathrm{op}}
 &\le M\ell+\tau+\tau B_{\mathrm{near}}.
\end{align*}
For the second inequality, Parseval gives squared column norm
$M s_2-\tau^2$ for $J$, and entrywise masking only decreases it.
For the third, $\|J\|\le M\ell+\tau$, while the removed matrix has
maximum absolute row sum at most $\tau B_{\mathrm{near}}$.
The entropy bound on $B_{\mathrm{near}}$ is the elementary binomial
tail bound on distances to a fixed cube point.

Let $R$ retain each of the $M$ coordinates independently with
probability $p=2n/M$. The zero-diagonal case of Tropp's random
principal-submatrix theorem \cite[Theorem 1.1]{tropp2008}, at moment
order two and $q=2\log M$, states
\[
 (\mathbb E\|RJR\|^2)^{1/2}
 \le C\left[q(\mathbb E\|RJR\|_{\max}^2)^{1/2}
       +\sqrt{pq}(\mathbb E\|JR\|_{1\to2}^2)^{1/2}
       +p\|J\|\right].
\]
Apply it to $J_{\mathrm{far}}$ and substitute the preceding bounds:
\begin{equation}\label{eq:bernoulli-gram}
 (\mathbb E\|R J_{\mathrm{far}}R\|^2)^{1/2}
 \le C[q\rho_r+\sqrt{nq s_2}+n\ell
               +(n/M)\tau(1+B_{\mathrm{near}})].
\end{equation}
The logarithm here is that of the original dimension $M$, not the
sample size.

To pass to a uniform subset of size $n$, use a uniform permutation of
all cube points, independently of $B\sim\operatorname{Bin}(M,2n/M)$.
Its first $B$ elements form the Bernoulli subset. On $B\ge n$, the
first $n$ principal submatrix is a compression of the larger one, so
its operator norm cannot be larger. The size-$n$ event is independent
of $B$, and $\mathbb P(B\ge n)\ge1/2$ for large $n$. Thus Markov's
inequality applied to \eqref{eq:bernoulli-gram} gives the same fixed-size
probability bound with a constant factor loss.

Conditional on no collisions, an independent size-$n$ sample is a
uniform size-$n$ subset, up to an irrelevant ordering. Collisions have
probability at most $n(n-1)/(2M)$. The probability of any sampled pair
violating the distance condition is at most $n^2B_{\mathrm{near}}/M$.
On their complement, the sampled matrices $J$ and $J_{\mathrm{far}}$
agree. Adding these exceptional probabilities proves
\eqref{eq:gram-bound}.
\end{proof}


\section{Tail covariance estimates at sparse cutoffs}\label{sec:sparse-gram}
The original principal-submatrix estimate uses the logarithm of the full
cube cardinality. We retain that estimate here. A more precise bound on
entries near orthogonal input pairs makes it applicable even when the
spectral cutoff is a sufficiently large fixed integer.

\begin{lemma}[Near-orthogonal entries]\label{lem:near-orthogonal}
There are constants $a_0,\delta_0,c,C>0$ such that the following holds.
Put $D=d+1$, let $r$ be odd with $r+2\le\delta_0D$, and let $h$ be the
Hamming distance between the augmented inputs $(1,x),(1,y)$.
If $a=|D-2h|/D\le a_0$, then
\begin{equation}\label{eq:near-orthogonal}
 \frac{|Q_r(x,y)|}{\tau_r}
 \le [C(a+\sqrt{(r+2)/D})]^{r+2}+C d^\beta e^{-cd},
 \qquad Q_r=K_{>r}.
\end{equation}
\end{lemma}
\begin{proof}
For $1\le j<D$, the Krawtchouk generating function gives
\begin{equation}\label{eq:sparse-kraw}
 \frac{|\mathsf K_j(h;D)|}{\binom Dj}
 \le [e^{3/2}(a+\sqrt{j/D})]^j.
\end{equation}
Indeed, set $b=\min(h,D-h)$. On the circle $|z|=t$, the generating
function has modulus at most
\[
 (1+t^2)^b(1+t)^{D-2b}\le \exp(Dt^2/2+Dat).
\]
Cauchy's coefficient bound and $\binom Dj\ge(D/j)^j$, with
$t=j/[D(a+\sqrt{j/D})]$, prove \eqref{eq:sparse-kraw}: both $Dt^2$
and $Dat$ are at most $j$.

Choose $a_0,\delta_0$ so small that
$e^{3/2}(a_0+\sqrt{\delta_0})\le e^{-1}$ and $\delta_0<1/8$.
On $j\le\delta_0D$, the derivative of
$j\log[e^{3/2}(a+\sqrt{j/D})]$ is at most $-1+1/2<0$.
Consequently \eqref{eq:sparse-kraw} is bounded by its value at $r+2$
throughout $r<j\le\delta_0D$. The sum of the corresponding positive
spectral masses $\binom Dj a_j$ is at most $\tau_r$.
For $\delta_0D<j\le(1-\delta_0)D$, the central-distance estimate
\eqref{eq:kraw} bounds the normalized Krawtchouk polynomial by
$C\sqrt D e^{-cD}$; its polynomial factor can be absorbed by decreasing
$c$. For $j>(1-\delta_0)D$, choose an odd $b_D$ near $D/4$.
Monotonicity of $a_j$ and Lemma~\ref{lem:sharp} bound the remaining
mass by
\[
 a_{b_D}\sum_{i<\delta_0D}\binom Di
 \le \operatorname{poly}(D)
       e^{-D[\mathsf h(1/4)-\mathsf h(\delta_0)]},
\]
where $\mathsf h(t)=-t\log t-(1-t)\log(1-t)$.
The entropy gap is positive. Finally $\tau_r\ge cD^{-\beta}$ by
Lemma~\ref{lem:sharp}, proving \eqref{eq:near-orthogonal}.
\end{proof}

\begin{lemma}[Sparse adaptive tail Gram matrices]\label{lem:sparse-gram}
Fix $A\ge1$, a failure probability $\delta>0$, and an odd integer
$r_0\ge11$. There is $\alpha_0>0$ such that, for any fixed
$0<\alpha\le\alpha_0$, the following assertions hold uniformly as
$d\to\infty$. Suppose $r$ is odd,
\[
 r_0\le r\le\alpha d,\qquad
 A D_r\log(D_r+1)\le n<A D_{r+2}\log(D_{r+2}+1).
\]
Put $u=r+2$, $\rho=r/(d+1)$, and $H=\log(e/\rho)$.
For each such specified $n$, with probability at least $1-\delta$,
\begin{align}
 \|(Q_r)_{SS}\|_{\mathrm{op}}
 &\le C_{L,\delta}\left[1+\frac{\sqrt{nd\,s_{2,r}}}{\tau_r}
                   +\frac{n\ell_r}{\tau_r}\right]\tau_r
 \notag\\
 &\le C_{L,\delta}
       [1+\sqrt{A\rho^{-1}H}+AH]\tau_r,
       \label{eq:sparse-coarse-gram}\\
 \|(Q_u)_{SS}-\tau_uI_n\|_{\mathrm{op}}
 &\le\bigl[C_{L,\delta}
       (\sqrt{A\rho H}+A\rho^2H)+o(1)\bigr]\tau_r.
       \label{eq:sparse-fine-gram}
\end{align}
The $o(1)$ is uniform over the stated range. In particular, after
reducing $\alpha_0$ if necessary,
$\|(Q_u)_{SS}\|_{\mathrm{op}}\le2\tau_r$ on the same event for all
sufficiently large $d$. Deterministically,
\begin{align}
 0\le\tau_r-\tau_u&\le C_L\tau_r/r,
       \label{eq:gram-sparse-band-trace}\\
 n\ell_u/\tau_r&\le C_LA\rho^2H,
       \label{eq:sparse-fine-eigenvalue}
\end{align}
where $\ell_v=\lambda_{v+1}$. No simultaneous event over different
sample sizes is asserted.
\end{lemma}
\begin{proof}
The dimension and spectral estimates already proved give
\begin{align*}
 n&\le C A\binom D{r+2}rH,&
 \log(2n)&\le C_A rH,\\
 n\ell_r/\tau_r&\le C_LAH,&
 nd\,s_{2,r}/\tau_r^2&\le C_LA\rho^{-1}H.
\end{align*}
For $u=r+2$, the exact two-step binomial ratio gives
\begin{align*}
 n\ell_u/\tau_r&\le C_LA\rho^2H,&
 nd\,s_{2,u}/\tau_r^2&\le C_LA\rho H.
\end{align*}
Moreover $\tau_r-\tau_u=\binom D{r+2}a_{r+2}$, so the level and
trace estimates prove \eqref{eq:gram-sparse-band-trace}.

We apply the proof of Lemma~\ref{lem:gram-general} with a different
deterministic mask. Retain pairs satisfying
\[
 \left|\frac{1+x\cdot y}{d+1}\right|\le a_n,
 \qquad a_n=\sqrt{128\log(2n)/d}.
\]
By choosing $\alpha_0$ sufficiently small, uniformly
$\log(2n)/d\le C_A\alpha\log(e/\alpha)+o(1)$ is as small as needed,
so $a_n\le a_0$. For independent uniform inputs, Hoeffding's bound
(and $a_n\ge2/(d+1)$) gives exceptional-pair probability
\[
 b_n\le2\exp(-a_n^2d/8)=2(2n)^{-16}.
\]
The exceptional set is translation invariant on the cube. Its size
is therefore the same in every row, so the masked-matrix proof applies
with $B_{\mathrm{near}}/M=b_n$. This mask is deterministic for each
fixed pair $(d,n)$; it does not depend on the observed sample.
Lemma~\ref{lem:near-orthogonal} supplies the retained-entry bound
\begin{equation}\label{eq:sparse-coherence}
 \frac{\eta_v}{\tau_v}
 \le\left[C\sqrt{\frac{\log(2n)+v+2}{d}}\right]^{v+2}
       +C d^\beta e^{-cd},\qquad v\in\{r,u\}.
\end{equation}
Thus exactly the argument of Lemma~\ref{lem:gram-general}, still with
$q=2\log M=2d\log2$, gives for every $z>0$
\begin{align}
 &\mathbb P\{\|(Q_v)_{SS}-\tau_vI_n\|>z\tau_v\}\notag\\
 &\quad\le\frac C{z^2}
 \left[\frac{q\eta_v}{\tau_v}
       +\frac{\sqrt{nq s_{2,v}}}{\tau_v}
       +\frac{n\ell_v}{\tau_v}+n(M^{-1}+b_n)\right]^2
       +n^2(M^{-1}+b_n).
 \label{eq:sparse-masked-tropp}
\end{align}
The collision and exceptional-pair terms tend to zero uniformly:
$n\ge A D_{r_0}\log(D_{r_0}+1)\to\infty$, and $\alpha_0$ can be chosen
so that $n\le e^{c_1d}$ with $2c_1<\log2$.

The coherence term $q\eta_v/\tau_v$ also tends to zero uniformly.
For $r\le\sqrt d$, its bracket in \eqref{eq:sparse-coherence} is at
most $C_A d^{-1/4}\sqrt{\log d}$, and its exponent is at least
$r_0+2\ge13$. Hence it is bounded, after multiplication by $q$, by
\[
 C d[C_A d^{-1/4}\sqrt{\log d}]^{13}
       +C d^{1+\beta}e^{-cd}=o(1).
\]
For $r>\sqrt d$, reducing $\alpha_0$ makes the bracket at most $1/2$,
so the bound is $Cd\,2^{-\sqrt d}+Cd^{1+\beta}e^{-cd}=o(1)$.
Using \eqref{eq:sparse-masked-tropp} for $v=r,u$, with failure
probability at most $\delta/2$ for each, now proves
\eqref{eq:sparse-coarse-gram} and \eqref{eq:sparse-fine-gram}.
For the latter, multiply by $\tau_u/\tau_r\le1$ before using the
spectral bounds, so no lower comparison of the two traces is needed.
Finally $\rho H\to0$ as $\rho\to0$, uniformly for $0<\rho\le\alpha$;
choose $\alpha_0$ to make the fixed part of the fine deviation below
$1/2$, and absorb the uniform $o(1)$. Since $\tau_u\le\tau_r$, this
gives the asserted upper bound $2\tau_r$.
\end{proof}

\section{The complete low-degree head}\label{sec:head}
\subsection{Squared balance on the target coordinates}
\begin{lemma}[Active-coordinate squared balance]
\label{lem:positive-active}
If $P$ has total Walsh degree at most $r$, then
\begin{equation}
 \mathbb E(fP)_-^2
 \ge\left(\frac12-\sqrt{\frac rk}\right)\|P\|_2^2.
 \label{eq:positive-active-balance}
\end{equation}
For every fixed $\delta>0$, there is a universal $A_\delta$ such that
if $r\le k/64$ and $n\ge A_\delta D_r\log(D_r+1)$, then, with
probability at least $1-\delta$,
\begin{equation}
 \mathbb P_n(fP)_-^2\ge\frac14\|P\|_2^2
 \quad\text{for every polynomial $P$ of degree at most $r$}.
 \label{eq:positive-active-sampling}
\end{equation}
\end{lemma}
\begin{proof}
For $j\in U$ put $D_jP=(P-P\circ\mathrm{flip}_j)/2$.
Because $f$ changes sign under these flips, the inequality
$|F(s)-F(t)|\le(|s|+|t|)|s-t|$ for $F(t)=t|t|$ gives
\[
 |\mathbb E fP|P||\le2\|P\|_2\|D_jP\|_2.
\]
Averaging over $j\in U$ and using
$\sum_{j\in U}\|D_jP\|_2^2\le r\|P\|_2^2$ proves
\eqref{eq:positive-active-balance} from
$(fP)_-^2=(P^2-fP|P|)/2$.
The vector of all Walsh characters of degree at most $r$ is isotropic
and has squared length $D_r$ at every input. Under $r\le k/64$,
the population bound is at least $3\|P\|_2^2/8$.
Apply Lemma~\ref{lem:squared-sampling} at accuracy $1/8$ and use
homogeneity to obtain \eqref{eq:positive-active-sampling}.
\end{proof}


\subsection{Separating the first omitted pair of Fourier levels}
\begin{lemma}[The first paired shell]
\label{lem:positive-first-shell}
Fix $A\ge1$, $\delta>0$, and an odd $r_0\ge11$.
There is a sufficiently small $\alpha>0$ such that, uniformly for all
large $d$, the following holds. Suppose $r$ is odd and
\[
 r_0\le r\le\alpha d,\qquad
 A D_r\log(D_r+1)\le n<A D_{r+2}\log(D_{r+2}+1).
\]
Write $\rho=r/D$, $H=\log(e/\rho)$, and $\tau=\tau_r$.
Then, with probability at least $1-2\delta$,
\begin{equation}
 \|(Q_r)_{SS}\|_{\mathrm{op}}\le C_{L,A,\delta}H\tau.
 \label{eq:positive-improved-gram}
\end{equation}
Deterministically,
\begin{align}
 n\lambda_{r+1}&\le C_L A H\tau,
 \label{eq:positive-first-eigen}\\
 n\lambda_{r+3}&\le C_L A\rho^2H\tau.
 \label{eq:positive-second-eigen}
\end{align}
\end{lemma}
\begin{proof}
Let $B_r=K_{(r,r+2]}$, so $Q_r=B_r+Q_{r+2}$.
The two levels in $B_r$ have the common normalized eigenvalue
$\ell=\lambda_{r+1}=\lambda_{r+2}$ and total dimension
\[
 m=\binom d{r+1}+\binom d{r+2}=\binom D{r+2}.
\]
Its diagonal is $\Delta=m\ell=\tau_r-\tau_{r+2}$.
Choose the feature vector
$v(x)=\sqrt\ell(\chi_V(x))_{r<|V|\le r+2}\in\mathbb R^m$.
Then $\|v(x)\|_2^2=\Delta$ and
$\mathbb E vv^{\mathsf T}=\ell I_m$.
For $A_i=v(X_i)v(X_i)^{\mathsf T}-\ell I_m$,
\[
 \|A_i\|_{\mathrm{op}}\le\Delta,\qquad
 \left\|\sum_{i=1}^n\mathbb E A_i^2\right\|_{\mathrm{op}}
 \le n\ell\Delta.
\]
The self-adjoint matrix Bernstein inequality
\cite[Theorem~1.4]{tropp2012}, with $h=\log(2m/\delta)$, therefore gives,
with probability at least $1-\delta$,
\begin{equation}
 \|(B_r)_{SS}\|_{\mathrm{op}}
 \le n\ell+C\sqrt{n\ell\Delta h}+C\Delta h.
 \label{eq:positive-band-bernstein}
\end{equation}
Here the nonzero eigenvalues of the sample Gram matrix equal those of
$\sum_i v(X_i)v(X_i)^{\mathsf T}$.

The sharp spectrum and maximality interval imply
\[
 n\le C A m rH,\quad
 \ell\le\frac{C_L\tau}{rm},\quad
 \Delta\le\frac{C_L\tau}{r},\quad
 h\le C_\delta rH.
\]
Thus every term in \eqref{eq:positive-band-bernstein} is at most
$C_{L,A,\delta}\tau H$.
The fine-tail conclusion of Lemma~\ref{lem:sparse-gram} for the current
adaptive cutoff $r$ gives $\|(Q_{r+2})_{SS}\|\le2\tau$ with probability
at least $1-\delta$
after reducing $\alpha$ and increasing $d_0$.
The triangle inequality proves \eqref{eq:positive-improved-gram}.
The first eigenvalue estimate is among the displayed spectral bounds;
the two-step binomial ratio
$\binom D{r+2}/\binom D{r+4}\le C\rho^2$ proves the second.
No independence between the two sample events is used.
\end{proof}


\subsection{The complete head at an exact sample size}
\begin{proposition}[Head norm and population leakage]\label{prop:head}
Let $r$ be odd, $r\le k/64$, $\rho=r/(d+1)$,
$H=\log(e/\rho)$, and $n=n_r$ from~\eqref{eq:exact-size}.
Outside an event of arbitrarily small fixed probability, every minimizer
of~\eqref{eq:obstacle}, with $P=z_{\le r}$, satisfies
\begin{align}
 \|z\|_K^2&\le C n a^2d^\beta,\label{eq:budget}\\
 \|P\|_2^2&\le C a^2\rho^{-\beta}H,\label{eq:head-norm}\\
 \|z-P\|_2^2&\le C a^2\rho^{2-1/400-\beta}H.
 \label{eq:head-leakage}
\end{align}
The estimates hold simultaneously for all minimizers.
\end{proposition}
\begin{proof}
The interpolation comparator gives~\eqref{eq:budget}; at the displayed
sample size, its $d$ term is absorbed by $nd^\beta$.
The one-sided squared sampling theorem and the parity balance estimate
give
\[
 \frac14\|P\|_2^2\le\frac1n\sum_i(f(X_i)P(X_i))_-^2
 \le\frac1n\sum_i(z(X_i)-P(X_i))^2
 \le\frac{\|z\|_K^2}{n}\|(Q_r)_{SS}\|.
\]
Lemma~\ref{lem:positive-first-shell} bounds the last Gram matrix by
$CH\tau_r$, proving~\eqref{eq:head-norm}.
The first omitted pair has common eigenvalue
$\lambda_{r+1}=\lambda_{r+2}$, dimension
$\binom{d+1}{r+2}\asymp D_r\rho^{-2}$, and trace at most
$C\tau_r/r$. Therefore
\[
 n_r\lambda_{r+1}\le C\tau_r H\rho^{2-1/400}.
\]
Since $\|z-P\|_2^2\le\lambda_{r+1}\|z\|_K^2$, this proves
\eqref{eq:head-leakage}. The sample size belongs to the interval in
Lemma~\ref{lem:positive-first-shell}, because $1/400<2$.
\end{proof}

In particular the low-degree head inherits the sign of $z$ outside a
set of vanishing population measure, and its mass in $(-a/2,a/2)$ tends
to zero. The estimate~\eqref{eq:head-norm} concerns the entire head;
the much weaker bound $\lambda_0\|z\|_K^2$ would not suffice below.

\section{Uniform middle-band estimates by Gaussian perturbation}
\label{sec:middle-pac}
This section uses the exact sample size~\eqref{eq:exact-size}. Put
$\tau=\tau_r$, $\rho=r/(d+1)$, and $H=\log(e/\rho)$.
Let $B=Cna^2d^\beta$ be a deterministic upper energy bound. The constants
may depend on fixed depth, polynomial sparsity, and a fixed upper
multiple $v_*/r$. Throughout the applications, $\epsilon\ge\rho^e$
for a sufficiently small fixed $e>0$; the positive-power remainders are
asserted only in this range. Auxiliary targets have uniform marginal,
or marginal density bounded by a fixed constant relative to uniform.
Their source is uniform, and the coupling is candidate independent
until the finite-net extension at the end of the section.

\begin{lemma}[A rare-event Gaussian PAC bound]\label{lem:gaussian-pac}
Let $\mathcal H$ be finite dimensional and
$F_h(X)=\max_{j\le M}\|T_{X,j}h\|$ for linear maps into Hilbert spaces.
Suppose
\[
 \operatorname{tr}(T_{X,j}^*T_{X,j})\le\Delta,
 \qquad \|T_{X,j}^*T_{X,j}\|\le L,
 \qquad \Pp\{F_h>z\}\le\Lambda\|h\|^2/z^2.
\]
Choose a centered isotropic Gaussian prior with variance $\sigma^2$ such
that $\Pp_G\{F_G>z/4\}\le\varepsilon_0$ uniformly in $X$.
Put $p_*=16\Lambda B/z^2+\varepsilon_0$ and suppose $p_*\le e^{-8}$.
With probability at least $1-\delta$, simultaneously for all
$\|h\|^2\le B$,
\[
 \Pp_n\{F_h>z\}\le C\sqrt{p_*}
 +C\frac{B/(n\sigma^2)+\log(2\sqrt n/\delta)/n}
              {\log(1/p_*)}.
\]
The conclusion remains true for auxiliary-averaged losses, with the
original $n$ independent observations as the sampling units.
\end{lemma}
\begin{proof}
For the bounded loss $\ell_g=1_{\{F_g>z/2\}}$, the binary PAC-Bayes
inequality holds simultaneously for every posterior $Q$:
\[
 \operatorname{kl}(\E_Q\Pp_n\ell_g\,\|\,\E_Q\Pp\ell_g)
 \le\frac{\KL(Q\|\Pi)+\log(2\sqrt n/\delta)}n.
\]
For completeness, the binomial exponential-moment bound
$\E\exp[n\operatorname{kl}(\Pp_n\ell\|\Pp\ell)]\le2\sqrt n$,
Fubini, Markov, and the entropy variational inequality prove this;
joint convexity of binary relative entropy passes from individual
losses to $Q$. The same argument applies to bounded losses in $[0,1]$
by Bernoulli randomization and Jensen.
Choose $Q=N(h,\sigma^2I)$, so $\KL(Q\|\Pi)=\|h\|^2/(2\sigma^2)$.
The triangle inequality gives empirical posterior loss at least
$(1-\varepsilon_0)\Pp_n(F_h>z)$ and population posterior loss at most
$p_*$. If its empirical value $q$ exceeds $\sqrt{p_*}$, then
\[
 \operatorname{kl}(q\|p')\ge q\log(1/p_*)
                 -q\log(e/q)
 \ge\tfrac14q\log(1/p_*)\qquad(p'\le p_*).
\]
Otherwise $q\le\sqrt{p_*}$. These alternatives prove the result.
Pointwise Gaussian estimates can be averaged over auxiliary variables
before applying the same bounded-loss inequality.
\end{proof}

The needed perturbation probabilities follow from Gaussian norm
concentration. In particular,
\begin{equation}\label{eq:gaussian-noise-options}
 \sigma^2=\frac{z^2}{C\Delta}
 \ \Longrightarrow\ 
 \varepsilon_0\le M e^{-c\Delta/L},\qquad
 \sigma^2=\frac{z^2}{C\Delta D_0},\quad D_0\ge\log(2M)
 \ \Longrightarrow\ 
 \varepsilon_0\le e^{-cD_0}.
\end{equation}
Indeed the mean norm is at most $\sigma\sqrt\Delta$ and the squared
Lipschitz constant is at most $\sigma^2L$.

\subsection{Dyadic completed pairs}
Index omitted pairs by $j\ge1$, with pair $j$ containing degrees
$r+2j-1,r+2j$. Group $t\le j<2t$ for $t=1,2,4,\ldots$, stopping at
$v_*$. If $I_t$ denotes a block, its kernel trace $\Delta_t$ and largest
eigenvalue $\Lambda_t$ satisfy
\begin{equation}\label{eq:dyadic-middle-spectrum}
 \Delta_t\le C(t/r)\tau,
 \qquad B\Lambda_t/a^2
       \le CH\rho^{-\beta-1/400}(C\rho)^{2t}.
\end{equation}
These follow by summing the paired level masses and iterating their
two-step binomial ratio. Thus the first omitted pair is included.
Give block $t=2^s$ threshold $a\epsilon_t$, where
$\epsilon_t=c\epsilon/(s+1)^2$ and $\sum_t\epsilon_t\le\epsilon$.
The triangle inequality reduces a bad event for the sum to the union
of block bad events. All fields may be sample dependent.

\subsection{Nested angular fluctuations}
For centered restriction to a nuisance subcube, the feature covariance
in a block has trace at most $\Delta_t$ and norm at most $C\rho\Delta_t$.
To verify the norm bound, diagonalize on the free coordinates. Every
nonconstant restricted character contains some specified free coordinate;
the sum of original Fourier weights containing that coordinate is at
most $(v_*/d)\Delta_t$. This bounds each restricted eigenvalue.
For each fixed nuisance order, reverse-martingale Doob gives
\[
 \Pp_X\left\{\max_q\Var_{C_{\pi,q}(X)}h>z^2\right\}
 \le\|h\|_2^2/z^2\le\Lambda_t\|h\|_K^2/z^2.
\]
The same bound holds after averaging random orders.
Use the first choice in~\eqref{eq:gaussian-noise-options}, with at most
$d+1$ stages. The perturbation failure is $e^{-c/\rho}$, and
\eqref{eq:dyadic-middle-spectrum} gives
$\log(1/p_*)\ge c\min(tH,1/\rho)$.
Lemma~\ref{lem:gaussian-pac} and dyadic summation yield
\begin{equation}\label{eq:nested-angular-pac}
 \Pp_{n,\pi}\left\{\max_q\Var_{C_{\pi,q}(X)}J>
                                      a^2\epsilon^2\right\}
 \le C\epsilon^{-2}\left[
 \frac{\rho^{-\beta}\log^5(r+2)}{rH}
           +\rho^{1-\beta}\log^4(r+2)\right]+\rho^c,
\end{equation}
with negligible confidence terms at the sample sizes under consideration.
Replacing $\epsilon^2$ by $p_{\rm bin}\epsilon^2$ divides the displayed
complexity terms by $p_{\rm bin}$. This controls subcube angular
fluctuations; it asserts no spherical weak-type estimate.

\subsection{Local common offsets}
Let $Y$ differ from $X$ only on an active face of dimension $m_f$.
For $h\le S$, the feature map
$h_0\mapsto[h_0(X)-h_0(Y,W)]_{W\in S_h(V)}$ has squared
Hilbert--Schmidt norm at most
\[
 C\rho(m_f+S)\Delta_t.
\]
The bound follows by averaging character differences: a degree at most
$v_*$ support meets the changed coordinates with probability at most
$C\rho(m_f+S)$. For the population mean-square maximum, the crude union
bound gives
\[
 \E\max_{h\le S}\E_{W\in S_h(V)}|h_0(X)-h_0(Y,W)|^2
       \le C(S+1)\Lambda_t\|h_0\|_K^2.
\]
Only the source and target marginals are used. Choose
$D_0=C[tH+\log(2S+2)]$ in the second Gaussian option. If
$S\le\rho^{-\Theta}$ and $2-1/400-\beta-\Theta>0$, the rare-event
logarithm is at least $ctH$. Consequently
\begin{equation}\label{eq:local-middle-pac}
 \Pp_{n,Y}\left\{
 \max_{h\le S}\E_{W\in S_h(V)}|J(X)-J(Y,W)|^2>a^2\epsilon^2
 \right\}
 \le C\epsilon^{-2}\rho^{1-\beta}(m_f+S)\log^4(r+2)+\rho^c.
\end{equation}
This proves the mean-square version needed to preserve a global gap.

\subsection{Radial windows and steps}
Let $R_b^m(h)$ be the normalized degree-$b$ nuisance sphere multiplier.
Its difference identity is
\[
 R_b^m(h)-R_b^m(h+1)=\frac{2b}mR_{b-1}^{m-1}(h).
\]
The Krawtchouk decay estimate implies, for $q\ge128\ell$,
$|h-q/2|\le\sqrt{2q\ell}$, and $b\le v_*$,
\[
 |R_b^m(h)-R_b^m(\lfloor q/2\rfloor)|^2
 \le C\ell b/m.
\]
Indeed sum the difference identity across the window and use
$|R_{b-1}^{m-1}(h)|\le C\exp[-c(b-1)q/m]$; maximizing
$(b/m)^2q\ell\exp[-c(b-1)q/m]$ proves the claim, with $b=0,1$
handled directly. Window feature variance is therefore at most
$C\ell\rho\Delta_t$.
The dimension-free spherical-mean maximal theorem~\cite{HKS} gives
\[
 \left\|\max_h|S_hh_0|\right\|_2\le C\|h_0\|_2,
\]
so it bounds the population loss of radial windows. Gaussian union over
at most $d^2$ windows and the preceding dyadic argument give
\begin{equation}\label{eq:middle-window-pac}
 \Pp_n\left\{\max_{\substack{q\ge128\ell\\|h-q/2|\le\sqrt{2q\ell}}}|S_hJ-S_{\lfloor q/2\rfloor}J|
                                   >a\epsilon\right\}
 \le C\epsilon^{-2}\rho^{1-\beta}\ell\log^4(r+2)+\rho^c.
\end{equation}
This invokes the maximal theorem for ordinary means, not for angular
variances or for spherical averages of squared functions.

For permutation-averaged subcube means the multiplier is
$A_b(q)=\binom{m-q}{b}/\binom mb$, with
$A_b(q)-A_b(q+1)=bA_b(q)/(m-q)$. Its absolute value is at most
$Cb/m$ for $b\le v_*$. The step feature variance is thus at most
$C\rho^2\Delta_t$. Jensen and reverse Doob give the population
$L^2$ maximal bound for $\mu_q$. The same proof yields
\begin{equation}\label{eq:middle-binomial-step-pac}
 \Pp_n\{\max_q|\mu_qJ-\mu_{q+1}J|>a\epsilon\}
       \le C\epsilon^{-2}\rho^{2-\beta}\log^4(r+2)+\rho^c.
\end{equation}

\subsection{Fixed couplings and a finite net}
Every proof above remains valid when the target active point is drawn
from a fixed candidate-independent coupling whose source is uniform,
whose target marginal density is bounded by a fixed constant, and whose
active changes stay within the indicated face. The pointwise Gaussian
feature bounds are unchanged, while population bounds gain only the
marginal-density constant. For a finite set of $N_*$ such couplings,
apply each PAC bound with confidence $\delta/N_*$. The additional
complexity is $\log N_*/n$, divided by the same rare-event logarithm.
A data-dependent head coupling is handled by the finite pattern net in
Section~\ref{sec:balanced-contacts}; no covariance concentration event is
unioned over that net.

\section{Spherical and permutation covariance envelopes}
\label{sec:angular-envelopes}
The spherical angular estimate requires a genuine maximal covariance
bound. A maximal theorem for ordinary spherical means alone would not
control averages of the squared field.

\begin{lemma}[Intrinsic positive-matrix concentration]\label{lem:intrinsic-psd}
Let independent positive matrices $A_i$ satisfy $A_i\le LI$. Put
$\mu=\|\sum_i\E A_i\|$ and $T=\operatorname{tr}\sum_i\E A_i$.
With probability at least $1-\delta$,
\[
 \left\|\sum_iA_i\right\|
 \le C\left[\mu+L\log\left(1+\frac{T}{L\delta}\right)\right].
\]
\end{lemma}
\begin{proof}
The matrix Laplace inequality and
$e^{A_i/L}\le I+(e-1)A_i/L$ give, with $M=\sum_i\E A_i$,
\[
 \E\operatorname{tr}(e^{\sum_iA_i/L}-I)
 \le\operatorname{tr}(e^{(e-1)M/L}-I)
 \le(e-1)(T/L)e^{(e-1)\mu/L}.
\]
Since the left random trace is at least
$e^{\|\sum_iA_i\|/L}-1$, Markov proves the result. Subtracting the
identity before Markov removes the ambient dimension.
\end{proof}

\subsection{A homogeneous spherical envelope}
On an $m$-cube give every degree-$b$ character eigenvalue $\kappa$ and
write $\pi_b=\binom mb\kappa$. Let $C_h$ be the centered covariance of
its evaluation feature on the sphere of radius $h$ about the all-plus
point.
\begin{lemma}[Homogeneous envelope]\label{lem:homogeneous-envelope}
For $m\ge5$ there is a positive operator $E_b\ge C_h$ for every $h$,
with
\[
 \|E_b\|\le\frac{\pi_b}{m-1},\qquad
 \operatorname{tr}E_b\le(b+1)\pi_b.
\]
If $b\le c\sqrt m$, then
$\operatorname{tr}E_b\le C_c\sqrt{b+1}\,\pi_b$.
\end{lemma}
\begin{proof}
The degree-$b$ subset module decomposes multiplicity freely into
harmonic components indexed by $0\le j\le\min(b,m-b)$, of dimensions
$M_j=\binom mj-\binom m{j-1}$. One elementary construction starts with
the kernels of the subset down operator and lifts them by the up
operator; the relation $DU-UD=(m-2j)I$ on degree $j$ gives the direct
sum and the stated dimensions. Permutation-invariant covariances are
scalar on these components. Sphere centering removes $j=0$; every
remaining component has dimension at least $m-1$. Take the maximum
covariance eigenvalue over $h$ separately on each component. Since every
$C_h$ has trace at most $\pi_b$, this proves the first two bounds.

For the sharper trace, a representative in component $j$ is
\[
 \prod_{i=1}^j(x_{2i-1}-x_{2i})
 \sum_{\substack{|B|=b-j\\B\subseteq\{2j+1,\ldots,m\}}}\chi_B(x).
\]
Its uniform squared norm is $2^j\binom{m-2j}{b-j}$.
On sphere $h$ each difference pair must contain exactly one minus sign,
an event of probability $2^j(h)_j(m-h)_j/(m)_{2j}$.
Consequently the fraction of homogeneous trace in component $j$ is
\begin{equation}\label{eq:harmonic-mass}
 T_j(h)=M_j\frac{4^j(b)_j(m-b)_j(h)_j(m-h)_j}{(m)_{2j}^2}
       \left[R_{b-j}^{m-2j}(h-j)\right]^2.
\end{equation}
Here $(x)_j$ is a falling factorial; impossible sphere parameters give
zero. Thus $\operatorname{tr}E_b/\pi_b=\sum_{j\ge1}\max_hT_j(h)$.

Write normalized Krawtchouk as a polynomial in the population mean
$t=1-2h/n$. Its monic version $P_l=(n)_lR_l/n^l$ satisfies
\[
 P_{l+1}(t)=tP_l(t)-\frac{l(n-l+1)}{n^2}P_{l-1}(t).
\]
The associated Jacobi matrix has symmetric real roots of absolute value
at most $2\sqrt{l/n}$. Pairing opposite roots and bounding the leading
coefficient therefore gives
\[
 |R_l^n(t)|^2\le e^{Cl^2/n}(t^2+4l/n)^l\qquad(l\le n/2).
\]
Set $s=4(h/m)(1-h/m)$. Falling-factorial estimates in
\eqref{eq:harmonic-mass}, followed by this root bound, give uniformly
for $b\le c\sqrt m$
\[
 T_j(h)\le e^{Cb^2/m}\binom bj
               s^j(1-s+Cb/m)^{b-j}
 \le C_c\Pp\{\operatorname{Bin}(b,s/(1+Cb/m))=j\}.
\]
Stirling's bounds imply
$\sup_t\Pp\{\operatorname{Bin}(b,t)=j\}
 \le C\sqrt{(b+1)/[(j+1)(b-j+1)]}$.
Summing over $j$ proves the square-root trace bound.
\end{proof}

\subsection{The complete highpass covariance}
Let $q\ge q_0$ be odd, with $q_0$ a sufficiently large fixed constant, $q\le d/8$, and
$k\le d^{1-\eta}$. Put $m=d-k$ and $\tau=\tau_q$.
\begin{proposition}[Maximal spherical angular covariance]
\label{prop:spherical-envelope}
There is a deterministic positive operator $A_X$ on the RKHS of $Q_q$
such that, for every field $t$,
\[
 \max_h\Var_{w\in S_h(V)}t(U,w)\le\langle t,A_Xt\rangle,
 \quad\|A_X\|\le C\tau/d,\quad
 \operatorname{tr}A_X\le C\tau d^2.
\]
Its translation mean has norm at most $Cq^2\lambda_{q+1}$.
If $q\le c_0\sqrt{d-k}$ for a sufficiently small fixed $c_0$, this
improves to $Cq^{3/2}\lambda_{q+1}$. Hence in that range, with probability
$1-\delta$, simultaneously for all $t$,
\[
 \sum_i\max_h\Var_{S_h(V_i)}t(U_i,\cdot)
 \le C\left[nq^{3/2}\lambda_{q+1}
 +\frac{\tau}d\log(1+nd^3/\delta)\right]\|t\|_K^2.
\]
The same statement holds with an independent uniform active input
averaged inside every summand.
\end{proposition}
\begin{proof}
At a fixed active input $u$, collapse the active Fourier coordinates.
For each nuisance support $B$, $b=|B|$, the corresponding mutually
orthogonal feature vectors have squared norm
\[
 \kappa_b=\sum_a\binom ka\lambda_{a+b}1_{a+b>q},
 \qquad \pi_b=\binom mb\kappa_b.
\]
Apply Lemma~\ref{lem:homogeneous-envelope} separately at each nuisance
degree. Put $s=\lfloor q/4\rfloor$, $\varepsilon=\beta/2$, and for
$b\ge s$ choose
\[
 w_b=(b+1)^{-1-\varepsilon}+s^{\beta-\varepsilon}\pi_b,
 \quad W=\sum_{b\ge s}w_b\le Cs^{-\varepsilon}.
\]
Weighted Cauchy shows that
$A_{\rm hi}=W\bigoplus_{b\ge s}(E_b/w_b)$ dominates the variance of
the sum of these homogeneous pieces. Since
$\pi_b/w_b\le s^{\varepsilon-\beta}$, its norm is at most
$Cs^{-\beta}/m$ and its trace at most $Cs^{-\beta}m^2$.

Uniform sign translation diagonalizes its mean. Pulling back from the
collapsed features to an original support with active degree $a$ and
nuisance degree $b$ cancels $\kappa_b$: its mean diagonal is at most
\[
 Cs^{-\varepsilon}(b+1)^{2+\varepsilon}\lambda_{a+b}.
\]
The paired spectrum implies, for every fixed $w>0$,
\[
 \sup_{j>q}(j+1)^w\lambda_j
     \le C_w(q+1)^w\lambda_{q+1}.
\]
Up to degree $d/4$, this follows from the growing paired binomial
denominators; beyond that point monotonicity gives exponential
suppression. Thus the high-piece mean is at most $Cq^2\lambda_{q+1}$.
With the sharper homogeneous trace, the power $2+\varepsilon$ becomes
$3/2+\varepsilon$ for $b\le c\sqrt m$. The remaining larger degrees
are exponentially suppressed if $q\le c_0\sqrt m$ with $c_0$ small,
giving $Cq^{3/2}\lambda_{q+1}$.

The omitted low nuisance degrees have total trace
$\Delta_{\rm lo}\le\exp[-cq\log(d/k)]$: conditional on total degree
$j>q$, fewer than $q/4$ nuisance coordinates require more than $3j/4$
active coordinates, with probability at most $2^j(k/d)^{3j/4}$.
The equal-weight envelope
$s\bigoplus_{1\le b<s}E_b$ has norm at most
$Cq\Delta_{\rm lo}/d$, negligible relative to $\tau/d$, and mean at
most $Cq^2\lambda_{q+1}$, or $Cq^{3/2}\lambda_{q+1}$ in the dilute
range. Degree zero has no angular variance. Combining the two pieces
and applying Lemma~\ref{lem:intrinsic-psd} proves the result.
\end{proof}

\subsection{Permutation means}
For a translation-invariant spectral subkernel of trace $\tau$ and
largest normalized eigenvalue $\lambda$, and a nuisance permutation
$\pi$, let $L_{X,\pi,q}$ be the feature
averaged on the first-$q$ free-coordinate cell and
$\overline L_{X,q}=\E_\pi L_{X,\pi,q}$. Put $J_d=1+\lceil\log_2(d+1)\rceil$.
There is a positive envelope $A_X^{\rm perm}$ satisfying
\begin{align*}
 \E_\pi\max_q|\langle h,L_{X,\pi,q}-\overline L_{X,q}\rangle|^2
       &\le\langle h,A_X^{\rm perm}h\rangle,\\
 \|A_X^{\rm perm}\|&\le CJ_d^2\tau/d,\qquad
 \operatorname{tr}A_X^{\rm perm}\le CJ_d^2\tau,\\
 \|\E A_X^{\rm perm}\|&\le CJ_d^2\lambda.
\end{align*}
To construct it, express each prefix as a sum of at most $J_d$ disjoint
dyadic intervals and apply Cauchy. On a fixed dyadic scale, interval
feature differences select disjoint Walsh support families, so total
trace is at most $\tau$. Centering over permutations removes the
trivial representation. The one-dimensional sign representation is
absent: a dyadic interval feature depends only on an ordered partition
into three coordinate blocks; for $m>3$ its stabilizer contains a
transposition. Every remaining irreducible component has dimension
at least $m-1$, giving the norm gain.
Translation averaging is diagonal; for each support the disjoint
interval indicators have total variance at most one per scale, giving
the mean bound.

For source inputs, Lemma~\ref{lem:intrinsic-psd} now gives a uniform
empirical bound of order
$J_d^2[n\lambda+\tau\log(nd/\delta)/d]$.
For a contact selected from the finite head-pattern net, use Gaussian
PAC rather than unioning this matrix event. In a dyadic band the envelope
has trace $CJ_d^2\Delta_t$, norm smaller by $d$, and population mean
$CJ_d^2\Lambda_t$. Gaussian perturbation with variance
$z^2/(CJ_d^2\Delta_t)$ has uniform noise failure $e^{-cd}$.
Lemma~\ref{lem:gaussian-pac} and dyadic summation therefore give
\begin{align*}
 &\Pp_{n,\rm contact}\left\{
 (\E_\pi\max_q|C_{\pi,q}J-\mu_qJ|^2)^{1/2}>a\epsilon\right\}\\
 &\qquad\le C\epsilon^{-2}\frac{\rho^{-\beta}\log^2d\log^5(r+2)}{rH}
 +\rho^c+\frac{C\log N_*\log(r+2)}{nH}.
\end{align*}
This is a positive power of $\rho$ in the large-$r$ branch
$r\ge\rho^{-4/5}$, for the small threshold powers used below.

\subsection{Coupled contacts in the small-cutoff branch}
The source event in Proposition~\ref{prop:spherical-envelope} cannot be
unioned over the head-pattern net. Instead compare a contact $U'$ with
the original active input $U$. For a middle-band block $I_t$, the feature
map
\[
 h_0\longmapsto[h_0(U',W)-h_0(U,W)]_{W\in S_h(V)}
\]
has trace at most $C\rho m_f\Delta_t$, uniformly at every radius:
only the at most $m_f$ changed active coordinates enter the character
difference. The full mean-square maximal difference has population
bound
\[
 \E\max_h\E_{W\in S_h(V)}|h_0(U',W)-h_0(U,W)|^2
             \le Cr^{3/2}\Lambda_t\|h_0\|_K^2.
\]
Indeed bound each endpoint's second moment by its variance plus squared
mean, then use the dilute envelope and the ordinary HKS maximal theorem.
The contact's bounded marginal density changes only the constant.
Gaussian norm concentration, a union over at most $d+1$ radii, and
Lemma~\ref{lem:gaussian-pac} therefore give
\begin{align*}
 &\Pp_{n,\rm contact}\left\{
 \max_h\E_{S_h(V)}|J(U',W)-J(U,W)|^2
                   >a^2\epsilon^2p_{\rm bin}\right\}\\
 &\qquad\le C\epsilon^{-2}\frac{\rho^{1-\beta}m_f\log^4(r+2)}{p_{\rm bin}}
          +\rho^c+\frac{C\log N_*\operatorname{polylog}(r)}n.
\end{align*}
Here $r\le\rho^{-4/5}$ and
$p_{\rm bin}=\rho^{\beta/2+30e}$, $\epsilon=\rho^e$.
The first block's population exponent is at least
$2-1/400-3\beta/2-6/5-32e>0$; the complexity exponent is at least
$1-3\beta/2-32e>0$.
Finally,
\[
 \Var_hJ(U',\cdot)\le2\Var_hJ(U,\cdot)
       +2\E_{S_h(V)}|J(U',W)-J(U,W)|^2.
\]
The source full-tail envelope, together with the uniform fine-tail
angular event, bounds the first term. The fine-tail event also holds
at the contact. This supplies the required full-tail spherical angular
control at adaptive balanced contacts without an invalid matrix union.

\section{Radial feature bounds and endpoint anchors}
\label{sec:radial-geometry}
Throughout this section the nuisance dimension $m=d-k$ is comparable
to $d$. The bounds also hold after a positive spectral restriction.

\begin{lemma}[Uniform Krawtchouk decay]\label{lem:uniform-kraw}
For normalized Krawtchouk polynomials,
\[
 |R_s^n(h)|\le
 \exp\left[-c\frac{\min(s,n-s)\min(h,n-h)}n\right].
\]
\end{lemma}
\begin{proof}
Reflection of either index preserves absolute value, and normalized
duality gives $R_s(h)=R_h(s)$. It suffices to treat $R_b(a)$ with
$1\le a\le b\le n/2$. Pair $q=\max(1,\lfloor a/4\rfloor)$ negative
population signs with positive signs. In the average over uniform
$b$-subsets, any subset splitting a pair cancels with the subset obtained
by switching membership within its first split pair. Hence the absolute
average is at most the probability that no pair splits.
For $a\ge4$, expose pairs successively. Conditional on previous
non-splitting patterns, at least $b/2$ selected and $n/4$ unselected
positions remain. Each next pair splits with probability at least
$b/(4n)$. Since $q\ge a/7$, the bound follows. For $a=1,2,3$, a single
pair splits with probability at least $b/n$, which gives the same bound
after reducing the absolute constant.
\end{proof}

\begin{lemma}[Sphere increments and windows]\label{lem:radial-increments}
Let $L_h$ be the kernel feature averaged on nuisance sphere $h$.
Then, uniformly for $0\le h<m$,
\[
 \|L_h-L_{h+1}\|_K^2\le Cd^{-\beta}
        [1+\min(h,m-1-h)]^{\beta-2}.
\]
Consequently, if $q\ge128\ell$, $q\le m$, and
$|h-q/2|\le\sqrt{2q\ell}$, then
\[
 \|L_h-L_{\lfloor q/2\rfloor}\|_K^2
          \le Cd^{-\beta}q^{\beta-1}\ell.
\]
For an active modification of distance $l$,
\[
 \|K_{(u,v)}-L_h(u',v)\|_K^2
              \le Cd^{-\beta}(l+h+1)^\beta.
\]
\end{lemma}
\begin{proof}
Let $\pi_j=\binom mj\kappa_j$ be the nuisance spectral distribution
obtained by fixing the active input. Its degree is no greater than the
full spectral degree, so $\sum_{j\ge t}\pi_j\le C(t+1)^{-\beta}$.
Tail integration gives
$\sum_j\pi_jj^2e^{-tj}\le Ct^{\beta-2}$ for $0<t\le1$.
Monotonicity of $\kappa_j$ also gives
$\sum_{j>3m/4}\pi_j\le2e^{-cm}$.
The exact difference formula yields
\[
 \|L_h-L_{h+1}\|_K^2
 =\frac4{m^2}\sum_{j=1}^m\pi_jj^2
                      [R_{j-1}^{m-1}(h)]^2.
\]
By reflection assume $h\le(m-1)/2$. For $j\le m/2$, apply
Lemma~\ref{lem:uniform-kraw} and the exponential second-moment bound.
For $m/2<j\le3m/4$, the Krawtchouk factor is at most $Ce^{-ch}$ and
the spectral mass at most $Cm^{-\beta}$. The remaining degree range
is exponentially small. This proves the first claim, including $h=0$
by the unweighted bound $\E j^2\le Cm^{2-\beta}$.
Telescoping feature norms across the central window proves the second.
The third follows from the endpoint cusp bound
$1-K(x,y)\le C(\dist(x,y)/d)^\beta$ and Jensen.
\end{proof}

\subsection{Midpoint and full-nuisance means}
For the total-degree highpass $Q_r$, the squared feature norm of nuisance
midpoint averaging is at most
\begin{equation}\label{eq:midpoint-trace}
 \Delta_{\rm mid}\le
 C\tau_r(k/d)^{r+1}+C\tau_r/d^2+C\sqrt d\,2^{-(d-k)}.
\end{equation}
To verify this, expand by nuisance degree $b$. At the middle sphere,
$|R_b^m(\lfloor m/2\rfloor)|\le C/m$ for $1\le b<m$.
For even $m=2t$ this follows directly from
$(1-z)^t(1+z)^t=(1-z^2)^t$; for odd $m=2t+1$ the generating function
is $(1+z)(1-z^2)^t$. Divide coefficients by $\binom mb$ and use
reflection to bound the remaining ratios by their first internal value.
The $b=0$ mass requires an entirely active support of degree above $r$,
giving the first term. The $b=m$ mass is exponentially small by spectral
monotonicity. Full nuisance averaging retains only $b=0$ and therefore
has just the first term.

Uniform translation makes the population endpoint operator norm at most
$\lambda_{r+1}$. For source inputs, Lemma~\ref{lem:intrinsic-psd}
applies directly. For a fixed contact coupling with bounded target
marginal, use a scalar Gaussian prior of variance
$a^2\epsilon^2/(C\Delta_{\rm mid}H)$ instead. Its uniform perturbation
failure is $e^{-cH}$; the shifted KL divided by $n$ is at most
$C\epsilon^{-2}d^\beta\Delta_{\rm mid}H$.
Lemma~\ref{lem:gaussian-pac} gives coupled endpoint bad fraction
\[
 \rho^c+C\epsilon^{-2}d^\beta\Delta_{\rm mid}
       +\frac{C[\log N_*+\log(n/\delta)]}{nH}
\]
simultaneously for a finite net of $N_*$ couplings. Polynomial sparsity
makes the pure-active term negligible for growing $r$; the $d^{-2}$
term is harmless. This avoids a union of endpoint matrix concentration
events over head-dependent contacts.

\section{Head recurrence and local drift}\label{sec:head-recurrence}

\begin{lemma}[Uniform recurrence of a head window]\label{lem:head-recurrence}
Let $P$ range over a $D_0$-dimensional linear space. Partition a bounded
range into $J_0$ fixed intervals. For a center and an independent nuisance
order, call it bad if its interval has relative mass less than $q$ in
some nested nuisance cell. For every $p_0>0$, the empirical bad fraction
is, simultaneously for all $P$, at most
\[
 Cq/p_0+CJ_0p_0+C(d+1)e^{-Lq/16}
 +C\sqrt{\frac{D_0\log(2+D_0J_0(d+1)(L+1))
                         +\log(1/\delta)}n},
\]
with probability $1-\delta$, provided $Lq\ge10$.
The loss may be averaged over nuisance orders and auxiliary active
inputs. For canonical head-dependent face contacts, the same bound has
an additional $O(D_0m_f)$ complexity and a fixed marginal-density factor.
\end{lemma}
\begin{proof}
For a fixed head and fixed active input, ignore intervals of population
mass less than $p_0$, losing at most $J_0p_0$. If $A$ is another interval,
under the law conditioned on membership in $A$, the reciprocal of its
conditional likelihood is a nonnegative supermartingale. Ville's
inequality shows that its cell mass drops below $q$ on at most a
$q/\Pp(A)$ fraction of this conditional law. Summing gives the fixed-head
bound $q/p_0+J_0p_0$.

For uniformity, draw $L$ independent ghost points in every nested cell.
Use the event that some cell has at most $3Lq/2$ ghost points in the
center's interval. A truly bad center triggers this ghost event with
probability at least $1/2$. If all cell masses exceed $2q$, the ghost
event has probability at most $(d+1)e^{-Lq/16}$ by Chernoff. Its fixed-head
expectation is therefore bounded by $2q/p_0+J_0p_0$ plus this term.
For $N$ augmented observations, all ghost labels use only
$O(NJ_0(d+1)(L+1))$ affine comparisons in the head coefficients.
The arrangement bound gives VC dimension
$CD_0\log(2+D_0J_0(d+1)(L+1))$. VC entropy, symmetrization, and Jensen
for ghost averaging prove the stated uniform bound.

For a canonical contact, first fix its face sign/order pattern.
The chosen contact is then fixed for each auxiliary random seed and
all ghost head evaluations remain linear in the coefficients.
The additional $O(4^{m_f})$ pattern comparisons add $O(D_0m_f)$ to
the VC dimension. Population bad mass grows by at most the contact
marginal-density constant. This is a bounded head-only argument; it
makes no covariance-event union over contact patterns.
\end{proof}

\begin{lemma}[Maximal local head drift]\label{lem:head-drift}
For degree at most $r$ and $S\le(d-k)/2$, the population local drift is
bounded by
\[
 \E_X\max_{h\le S}\E_{W\in S_h(V)}|P(U,W)-P(U,V)|^2
 \le C\frac{rS}{d-k}\log^2(2S)\|P\|_2^2.
\]
The same bound holds empirically, uniformly over the degree-$r$ space,
up to a fixed constant when $n\ge CD_r\log(D_r/\delta)$.
For head-dependent balanced contacts, the threshold bad fraction is
controlled by this population bound and bounded VC transfer.
\end{lemma}
\begin{proof}
Couple the spheres by a uniformly random distinct-coordinate nuisance
path. Every path prefix is a union of at most $O(\log S)$ dyadic blocks.
Cauchy bounds the maximum squared increment by $O(\log S)$ times the
sum of squared dyadic block increments. A block of length $l$ changes
a degree at most $r$ character with probability at most $rl/(d-k)$.
Translation orthogonality therefore bounds the expected block quadratic
form by $Clr/(d-k)$ times the identity. On each dyadic scale the total
length is at most $2S$, proving the population estimate.
The resulting positive envelope has mean norm
$\mu=C(rS/(d-k))\log^2(2S)$ and trace at most $D_r\mu$.
The trace bound is pointwise (translation changes character signs), so
the envelope norm is at most $D_r\mu$. Matrix Chernoff gives its
empirical bound for the stated sample size.
For a balanced contact, the population bound gains only its fixed target
marginal density. The indicator of a threshold failure is determined,
on a face pattern cell, by at most $S+1$ quadratic forms in the head
coefficients. Its VC dimension is $CD_r[m_f+\log(2S)+1]$, so bounded
statistical transfer applies. This avoids multiplying a matrix bound
by the number of head patterns.
\end{proof}

The head cap used in recurrence is justified similarly by
$\|P\|_2^2\le Ca^2\rho^{-\beta}H$ and bounded threshold VC transfer.
Taking cap $a\rho^{-\beta/2-10e}$, interval width at most $a/16$,
$p_0=\rho^{\beta/2+20e}$, and $q=\rho^{\beta/2+30e}$ makes the
uniform recurrence failure a positive power of $\rho$ at the exact
sample sizes. Markov over the auxiliary nuisance order then gives
simultaneous recurrence along at least three quarters of orders outside
a positive-power exceptional set of center/contact pairs.

\section{Transport of the all-input margin}\label{sec:transport}
Write an input as $(u,v)$, with $u\in\{-1,1\}^k$ active and
$v\in\{-1,1\}^{m}$ nuisance, where $m=d-k$. For a nuisance function, let $S_h$
denote averaging on the Hamming sphere of radius $h$ about its indicated
center. For a permutation $\pi$ of the nuisance coordinates, let
$C_{\pi,q}$ average over the subcube obtained by freeing the first $q$
coordinates, and put $\mu_q=\E_\pi C_{\pi,q}$.
Choose an odd $v_*=C_0r+O(1)$ and split
\[
 z=P+J+T,\qquad P=z_{\le r},\quad J=z_{(r,v_*]},\quad T=z_{>v_*}.
\]
The letter $v_*$ denotes a cutoff, not a nuisance input.

\begin{lemma}[Effective profiles]\label{lem:effective-profile}
Fix $\eta>0$ and the exact sample sizes of Theorem~\ref{thm:main}.
There are constants $g,\Theta>0$ such that the following holds with
probability tending to one, simultaneously for all globally separated,
training-compatible fields satisfying~\eqref{eq:budget} and
\eqref{eq:head-norm}. Put $S=\lfloor\rho^{-\Theta}\rfloor$.
Choose a random active face of dimension $m_f$, where $m_f$ is bounded
by a fixed power of $H$, through the observed active input $U_i$.
For a uniform auxiliary point, or a canonical balanced contact $u$
from Section~\ref{sec:balanced-contacts} on an eligible face, define
\[
 W_{i,u,h}=P(u,V_i)+J(X_i)+S_hT(u,V_i),\qquad0\le h\le2S.
\]
Outside an empirical fraction at most $\rho^g$ of center--face--point
pairs, the profile has margin $a(1-o(1))$ at every radius, and
\[
 \sign W_{i,u,S}=\sign P(u,V_i).
\]
The same margin conclusion holds at the distinguished input $u=U_i$
outside a fraction $\rho^g$ of centers; at radius zero its value equals
$z(X_i)$ exactly. The statement concerns auxiliary averages, rather
than a maximum over arbitrary active points.
\end{lemma}

\subsection{Uniform high-frequency witnesses}
\begin{lemma}[Fine-tail evaluation and contrast Gram matrices]
\label{lem:fine-cell-gram}
Let $v_*=C_0r+O(1)$, with $C_0$ a sufficiently large fixed constant.
Assume $\log n=O(rH)$, $k=O(rH)$, and $\log d=O(H)$.
At every center allow all active translations and a polynomial number
of nuisance sphere averages, adjacent increments, and radial window
contrasts. Give each scalar contrast a deterministic variance upper
bound at least an inverse polynomial in $d$, and normalize by that bound.
With probability $1-\exp(-crH)$, every adaptive choice of one such
normalized contrast per center has Gram norm at most a fixed constant.
Moreover,
\[
 \sum_i\sup_{u,h}\Var_{w\in S_h(V_i)}T(u,w)
       \le C\frac{\tau_{v_*}}d\|T\|_K^2.
\]
\end{lemma}
\begin{proof}
For translated, uniformly bounded signed evaluation averages, Fourier
orthogonality bounds the second moment of a cross-center covariance by
$C\lambda_{v_*+1}\tau_{v_*}$. The analogous Hilbert--Schmidt bound holds
for the centered spherical evaluation operators. Their individual
operator norms are at most $C\tau_{v_*}/(d-k-1)$: nuisance permutation
invariance removes the trivial representation, and each remaining
irreducible component has dimension at least $d-k-1$.
A union over the at most $2^k\operatorname{poly}(d)$ choices at each
center, followed by Markov and Gershgorin, bounds all normalized
cross-entries by $1/n$. Its failure probability is at most
\[
 Cn^4 4^k\operatorname{poly}(d)
                 \lambda_{v_*+1}\tau_{v_*}.
\]
The spectral entropy bound makes this at most $\exp(-crH)$ after
increasing $C_0$. The block Hilbert--Schmidt argument gives the second
assertion. All choices may be made after observing the field.
\end{proof}

\subsection{Proof of the effective-profile lemma}
\begin{proof}[Proof of Lemma~\ref{lem:effective-profile}]
Choose $e=1/40000$, $\epsilon=\rho^e$,
$p_{\rm bin}=\rho^{\beta/2+30e}$, and
$S=\lfloor\rho^{-(1-\beta)/2}\rfloor$.
Lemma~\ref{lem:head-drift}, the head norm bound, and
\eqref{eq:local-middle-pac} give, outside positive-power fractions of
source/contact pairs, simultaneously for $h\le2S$,
\[
 \E_{S_h(V_i)}|P(u,W)-P(u,V_i)|^2\le a^2\epsilon^2,
 \quad
 \E_{S_h(V_i)}|J(u,W)-J(X_i)|^2\le a^2\epsilon^2.
\]
The fine-tail angular Gram bound gives the same smallness for the
fluctuation of $T$ about its sphere mean. On each sphere Chebyshev
leaves a point where all deviations are $O(a\epsilon)$; the all-input
margin of $z$ therefore gives $|W_{i,u,h}|\ge a(1-O(\epsilon))$.
At the observed input and radius zero the decomposition is exact.

We now anchor the sign. Lemma~\ref{lem:head-recurrence} supplies, for
at least three quarters of nuisance orders, a head-value window of
width at most $a/16$, containing $P_0=P(u,V_i)$, with mass at least
$p_{\rm bin}$ in every nested cell. The statement holds uniformly for
the original head and for its canonical contacts. The cap failure,
ghost recurrence error, and contact-pattern mismatch all have a positive
power bound at the exact sample sizes.

Suppose first that $r\le\rho^{-4/5}$. The dilute envelope and its
contact extension give, outside a positive-power exceptional fraction,
\[
 \max_h\Var_{S_h(V_i)}(J+T)(u,\cdot)
                  \le a^2p_{\rm bin}\epsilon^2/100.
\]
Indeed after multiplying the envelope by the energy budget and dividing
by $p_{\rm bin}\epsilon^2$, its coherent term has exponent at least
$2-1/400-3\beta/2-6/5-32e>0$; its diagonal term has a larger exponent.
Average the recurrent window mass over nuisance permutations. A uniformly
freed $q$-subcube is exactly a mixture of spheres with radius
$Z\sim\operatorname{Bin}(q,1/2)$. Put
$\ell=\log(100/p_{\rm bin})$. At least a fixed multiple of $p_{\rm bin}$
of its joint mass lies in the matching head window and in
$|Z-q/2|\le\sqrt{2q\ell}$. Angular Chebyshev cannot remove this entire
mass. At a surviving point the global margin, the head-window width,
and the middle/fine radial-window bounds yield
\[
 |P_0+S_{\lfloor q/2\rfloor}(J+T)(u,V_i)|\ge3a/4
                 \qquad(2S\le q\le m).
\]
The radial-window estimate \eqref{eq:middle-window-pac} controls $J$;
Lemma~\ref{lem:radial-increments} and the uniform fine-tail Gram event
control $T$. Their fine-tail bad fraction is at most
$C\epsilon^{-2}S^{\beta-1}\ell$, with no division by $p_{\rm bin}$.
Adjacent radial steps are small, so these gapped expressions cannot
change sign. The midpoint endpoint is small by
\eqref{eq:midpoint-trace} and its scalar-PAC contact extension.
The common sign is therefore $\sign P_0$. At radius $S$, the local
middle-band approximation transfers this sign to $W_{i,u,S}$.

For $r\ge\rho^{-4/5}$ use the hybrid profile
\[
 H_q=P_0+\mu_qJ(u,V_i)+S_{\lfloor q/2\rfloor}T(u,V_i),
                        \qquad2S\le q\le m.
\]
Equation~\eqref{eq:nested-angular-pac}, at threshold
$a^2p_{\rm bin}\epsilon^2$, controls the middle-band subcube variance.
Its leading exponents are at least
$4/5-3\beta/2-32e$ and $1-3\beta/2-32e$, both positive.
The permutation-mean estimate controls the difference from $\mu_qJ$.
For source inputs use its intrinsic matrix bound; for balanced contacts
use its Gaussian-PAC finite-net form. Markov over orders makes both
properties simultaneous along at least three quarters of orders.

Their intersection with the recurrent-window orders has mass at least
one half. In each good cell delete the middle deviation set, costing
at most $p_{\rm bin}/100$. Averaging the remaining matching mass over
orders leaves at least $0.49p_{\rm bin}$ in the binomial sphere mixture.
The binomial tail outside the central window costs at most
$p_{\rm bin}/100$. The fine-tail spherical angular event, which is
uniform over all active contacts, removes at most another
$p_{\rm bin}/100$. At a surviving point all three approximations hold.
The fine radial-window bound thus gives $|H_q|\ge3a/4$ for every $q$.
This step uses subcube angular control only for $J$ and spherical angular
control only for $T$.

The binomial step estimate~\eqref{eq:middle-binomial-step-pac} and the
fine sphere-step Gram bound make consecutive $H_q$ differ by $o(a)$.
The full-nuisance mean of $J$ and the midpoint mean of $T$ are small
by the endpoint bounds, including their contact-PAC form. Hence every
$H_q$ has sign $\sign P_0$. Convex averaging of the local middle-band
mean-square estimate gives
$|\mu_{2S}J(u,V_i)-J(X_i)|\le a\epsilon$.
Therefore $W_{i,u,S}$ has the same sign.

All error estimates above are powers of $\rho$ times fixed logarithms.
The smallest branch slack at $\beta=1/2$ is
$0.05-1/400-32e>0$. Recurrence errors and pattern mismatches are also
positive powers; Markov over auxiliary orders and the finitely many
uses reduces these exponents only by fixed factors. Taking, for example,
$g=5e$ after all reductions is valid. This proves the lemma, uniformly
in the field, including balanced contacts.
\end{proof}

\subsection{The two-scale witness}
\begin{lemma}[A wrong head value yields a witness]\label{lem:coarse-witness}
Let a face have dimension $m_f$, and normalize its head by two lower
same-sign amplitude quantiles. Write $Q=(P-t)/s$, with $s\ge a/2$, and
let $F$ be the projection of $Q$ onto the halfline of $\sign P$ with
absolute margin one. Fix an amplitude cap $R$ with
$8(m_f+1)\le R\le S$. At a wrong center outside the quantile exception
set put $b=|F|$. Suppose a same-head-sign quantile contact has an effective
profile as in Lemma~\ref{lem:effective-profile}.
One adaptive fine-tail scalar witness then has squared value divided
by its variance upper bound at least
\[
 ca^2d^\beta
 \begin{cases}
 b^2/(m_f+b+1)^\beta,&1\le b\le R/4,\\
 R^{2-\beta},&b>R/4.
 \end{cases}
\]
\end{lemma}
\begin{proof}
The head normalization gives a center--contact contrast at least $cab$
whenever the contact profile has its original sign. If $b\le R/4$,
inspect radius $h_0=\max(1,\lceil b\rceil)$; otherwise inspect
$h_0=\lfloor R/2\rfloor$. If the original sign is already present,
use the center--contact contrast, whose variance is at most
$Cd^{-\beta}(m_f+h_0+1)^\beta$.
Otherwise the sign anchor at $S$ forces an adjacent sign-changing pair
at some radius $h\ge h_0$. Its jump has magnitude at least $ca$, while
its variance is at most $Cd^{-\beta}(h_0+1)^{\beta-2}$.
These two alternatives give the displayed bounds. The common scalar
offset cancels in both contrasts.
\end{proof}

The adjacent-shell variance estimate follows from the exact
Krawtchouk multiplier formula and the cusp exponent $\beta$:
\[
 \Var_K(S_h-S_{h+1})\le Cd^{-\beta}(h+1)^{\beta-2}.
\]
The center--contact estimate follows from the kernel increment bound
$1-K(x,y)\le C(\dist(x,y)/d)^\beta$ and Jensen.
Lemma~\ref{lem:fine-cell-gram} adds one chosen witness per center.
Consequently, for $p=2-\beta$, the empirical head-only losses satisfy
\begin{align}
 \Pp_{n,\rm face}[1_{\rm wrong}b^p1_{b\le R/4}]
   &\le C(m_f+1)^\beta+CR^p\delta_{\rm geom}+1,
       \label{eq:truncated-moment}\\
 \Pp_{n,\rm face}[1_{\rm wrong}1_{b>R/4}]
   &\le CR^{-p}+\delta_{\rm geom}.
       \label{eq:huge-set}
\end{align}
The geometric bad fraction obeys $\delta_{\rm geom}\le\rho^g$ after
quantile contact selection. The cap $R=\rho^{-t_{\rm rad}}$ is chosen
much smaller than the anchor horizon $S$. This separation ensures that
$R^p\delta_{\rm geom}\to0$.

\section{Balanced contacts and the sharper face obstruction}
\label{sec:balanced-contacts}

The maximum possible distance between two points on an active face is
its dimension. A more useful comparison preserves the distribution of
the center and gives the contact a bounded density. The resulting
distance has moments on the square-root scale. We first prove the
coupling statement, including the simultaneous tail bound needed below.

\begin{lemma}[A bounded-density coupling]\label{lem:balanced-coupling}
Let $\mu$ be uniform on $\{-1,1\}^m$, and let $\nu$ have density at
most $C\ge1$ relative to $\mu$. There is a single coupling
$X\sim\mu$, $Y\sim\nu$ such that, for every $t\ge0$,
\begin{equation}\label{eq:balanced-tail}
 \Pp\{\dist(X,Y)>2\sqrt{m(\log C+2t)}+t\}\le2e^{-t}.
\end{equation}
Consequently, for every $q\ge1$,
\[
 \|\dist(X,Y)\|_q\le C_0\sqrt{m(\log C+q)}+C_0q.
\]
Two measures of densities at most $C_1,C_2$ can likewise be coupled
with distance $L$ satisfying
\begin{equation}\label{eq:two-density-distance}
 \|L\|_q\le C_0\sqrt{m(\log C_1+\log C_2+q)}+C_0q.
\end{equation}
\end{lemma}
\begin{proof}
Expose coordinates in a fixed order and write
\[
 p_i(y_{<i})=\nu\{Y_i=1\mid Y_{<i}=y_{<i}\},\quad
 \delta_i=|p_i-1/2|,\quad S_i=\sum_{j\le i}\delta_j^2.
\]
On a zero-probability prefix the conditional probability may be defined
arbitrarily. Let $L_i$ be the likelihood ratio of the first $i$
coordinates under $\nu$ against their uniform law. Under the uniform
law, conditional on a positive-likelihood prefix,
\[
 \E[\sqrt{L_i}\mid Y_{<i}]
 =\sqrt{L_{i-1}}\,
       \frac{\sqrt{p_i}+\sqrt{1-p_i}}{\sqrt2}.
\]
The square of the last factor is
\[
 \frac{1+\sqrt{1-4\delta_i^2}}2
 \le1-\delta_i^2\le e^{-\delta_i^2}.
\]
Thus $\sqrt{L_i}e^{S_i/2}$ is a nonnegative supermartingale. Since
$L_m\le C$,
\begin{equation}\label{eq:conditional-bias-energy}
 \E_\nu e^{S_m/2}
 =\E_\mu L_me^{S_m/2}
 \le\sqrt C\,\E_\mu\sqrt{L_m}e^{S_m/2}\le\sqrt C.
\end{equation}

Given the full joint past, maximally couple a fresh uniform Bernoulli
$X_i$ with a Bernoulli variable $Y_i$ of parameter $p_i(Y_{<i})$.
The conditional marginal of $X_i$ is uniform even after conditioning
on the full past; hence $X$ is uniform, and $Y$ has law $\nu$.
The disagreement indicator has conditional mean $\delta_i$.
For $D=\dist(X,Y)$ and $A=\sum_i\delta_i\le\sqrt{mS_m}$,
the Bernoulli exponential supermartingale gives
\[
 \E\exp\{D-(e-1)A\}\le1.
\]
Because $e-1<2$, the event $D>2\sqrt{mS_m}+t$ has probability at most
$e^{-t}$. Equation~\eqref{eq:conditional-bias-energy} also gives
$\Pp\{S_m>\log C+2t\}\le e^{-t}$. Their union proves
\eqref{eq:balanced-tail}; integration proves the moment bound.
Finally, glue the two couplings through a common uniform intermediate
point and use the triangle inequality and Minkowski's inequality.
\end{proof}

Fix a sufficiently small positive constant $\zeta_0$.
Write $\mu$ for uniform probability on the face and set $\sign0=1$.
A face is eligible when each head-sign class
$S_\sigma=\{x:\sign P(x)=\sigma\}$ has mass at least $2\zeta_0$ and
$\mu\{|P|<a/2\}<\zeta_0/2$.
For $\sigma\in\{-1,1\}$ define the lower fixed-mass quantile
\[
 b_\sigma=\min\{b\ge0:\mu(S_\sigma\cap\{|P|\le b\})\ge\zeta_0\}.
\]
All points tied at the quantile value are included. The strict lower set
has mass less than $\zeta_0$, so ties cause no ambiguity in the projection
estimate below. Eligibility implies $b_\sigma\ge a/2$ and the usable
contact sets
\[
 A_\sigma=S_\sigma\cap\{a/2\le|P|\le b_\sigma\}
 \quad\hbox{satisfy}\quad\mu(A_\sigma)\ge\zeta_0/2.
\]
Put
\[
 s=(b_++b_-)/2\ge a/2,\qquad t=(b_+-b_-)/2,\qquad Q=(P-t)/s.
\]
Project $Q(x)$ onto the halfline $\{y:(\sign P(x))y\ge1\}$ and call
the resulting value $F(x)$. Then $|F|\ge1$ and $\sign F=\sign P$.
A displacement occurs only on the strict lower quantile set in one of
the two sign classes. There $Q\in[-1,1]$, so the displacement is at most
two and
\begin{equation}\label{eq:contact-projection-leakage}
 \|F-Q\|_2^2\le8\zeta_0.
\end{equation}
On ineligible faces all statistical losses are defined as zero.
Apply Lemma~\ref{lem:balanced-coupling} to the uniform laws on
$S_\sigma$ and $A_\sigma$, and mix the two couplings with weights $\mu(S_\sigma)$.
This defines a joint law of a uniform center $U$ and a contact $V$.
The contact has the same original head sign, its marginal density is
at most $2/\zeta_0$, and
\begin{equation}\label{eq:contact-distance-moments}
 \|\dist(U,V)+1\|_q
 \le C_{\zeta_0}\{\sqrt{mq}+q+1\}\qquad(q\ge1).
\end{equation}
Fix the coordinate order and all one-bit coupling conventions once.
The resulting contact law depends only on the sign and ordering
pattern of the head on the face. This observation permits uniform
control when the head itself depends on the training data.

\subsection{Uniformity for a head-dependent contact law}
\begin{lemma}[A finite net of face patterns]\label{lem:contact-pattern-net}
Let heads range over a $D_0$-dimensional real linear space.
Let $c_P(\mathcal F)$ record the signs, orderings, and fixed threshold
comparisons on an $m$-face that determine eligibility and the contact
sets. For $0<\epsilon<1/2$, there is a data-independent finite set of
heads $P_1,\ldots,P_M$ with
\begin{equation}\label{eq:contact-net-size}
 \log M\le C D_0\{m+\log(1/\epsilon)+1\}
\end{equation}
such that every $P$ has an index $j$ with
$\Pp_{\mathcal F}\{c_P(\mathcal F)\ne c_{P_j}(\mathcal F)\}\le\epsilon$.
With probability at least $1-\delta$ over $n$ independent uniform
training inputs, the same index has empirical disagreement, averaged
over faces through those inputs, at most
\begin{equation}\label{eq:contact-net-empirical}
 \epsilon+C\sqrt{\frac{D_0(m+1)+\log(1/\delta)}n}.
\end{equation}
\end{lemma}
\begin{proof}
Each face pattern is determined by at most $C4^m$ affine comparisons
in the head coefficients. On $N$ faces, the disagreement indicators
of a pair of heads are therefore determined by an arrangement of
$CN4^m$ hyperplanes in $2D_0$ parameters. The number of cells,
including cells on hyperplanes, is at most
$(CN4^m/D_0)^{CD_0}$ for $N\ge D_0$. Solving the shattering
inequality shows that the disagreement class has VC dimension at most
$CD_0(m+1)$; an extra $\log D_0$ is unnecessary here.

The VC epsilon-net theorem supplies a fixed auxiliary collection of
at most $C D_0(m+1)\epsilon^{-1}\log(2/\epsilon)$ faces meeting every
disagreement set of probability greater than $\epsilon$. Partition
coefficient space by the comparisons on those faces and choose a
representative in each nonempty cell. The arrangement bound gives
\eqref{eq:contact-net-size}. Two heads in one cell agree on the
auxiliary collection, so their population disagreement is at most
$\epsilon$. The auxiliary collection establishes existence of the net
and is independent of the training data.

The VC entropy bound gives uniform empirical deviation
$C\sqrt{(D_0(m+1)+\log(1/\delta))/n}$ for the disagreement class.
Jensen's inequality preserves this bound when the auxiliary face is
averaged before evaluating each training input. This proves
\eqref{eq:contact-net-empirical}.
\end{proof}

For each fixed net index, the contact law is an ordinary fixed
augmented-input distribution. Its source is uniform and its target
has bounded density. The Gaussian perturbation estimates used in
the middle-band arguments are uniform in the augmented observation,
and their population estimates use only these marginals.
Applying those arguments simultaneously to the net costs
$\log M$ in the PAC-Bayes numerator. On an agreeing face the entire
conditional contact law agrees, so the discarded fraction
\eqref{eq:contact-net-empirical} transfers the conclusion to the
data-dependent head. At the specified sample sizes, $D_0=D_r$,
$m\le CrH$, and $\epsilon=\rho^c$, both this fraction and
$\log M/n$ are positive powers of $\rho$, up to logarithmic factors.

The marginal statement alone would not justify this substitution:
a coupling can concentrate on a source/contact relation having small
product measure. The finite-pattern argument supplies the required
uniformity for that relation. Likewise, one must not union a
$D_r$-dimensional matrix concentration bound over all net indices.
Head-only drift and recurrence losses instead use their bounded VC
estimates; the middle-band extensions use Gaussian PAC-Bayes.

\subsection{Weighted witnesses and bounded statistical losses}
Write $\ell=\dist(U,V)$, and let $b=|F(U)|$ for the projected,
contact-normalized head. Use separate scales: the effective radial
profile is anchored at a horizon $S$, whereas amplitudes are capped
at $R\ll S$, with $m\ll R$. At a good wrong center outside the
contact-quantile exception, the proof of the two-scale witness lemma
with the actual distance $\ell$ gives a normalized scalar witness of
squared value at least
\begin{equation}\label{eq:distance-witness}
 ca^2d^\beta
 \begin{cases}
 b^2/(\ell+b+1)^\beta,&b\le R/4,\\
 R^{2-\beta},&b>R/4.
 \end{cases}
\end{equation}
For completeness, test the contact profile at $h_0=\lceil b\rceil$
in the first case and at $h_0=\lfloor R/2\rfloor$ in the second.
If it has the original head sign, compare it with the observed center.
Otherwise an adjacent sign change between $h_0$ and $S$ supplies the
witness. The contrast and increment variances are bounded respectively
by $Cd^{-\beta}(\ell+h+1)^\beta$ and
$Cd^{-\beta}(h+1)^{\beta-2}$. This proves
\eqref{eq:distance-witness} while allowing all intermediate sign relays.

Put $p=2-\beta$ and retain the full contact score
\[
 e(b,\ell)=\frac{b^2}{(\ell+b+1)^\beta}.
\]
For $b\ge1$ it satisfies $e(b,\ell)\le b^p$ and
\begin{equation}\label{eq:full-contact-minimum}
 e(b,\ell)\ge 2^{-\beta}
       \min\{b^2/(\ell+1)^\beta,b^p\}.
\end{equation}
Keeping both terms will improve the analytic obstruction.
The uniform fine-tail Gram event permits one adaptively chosen scalar
witness per training center. Sum that deterministic inequality and
then average over the face and the canonical contact law. The
comparator budget and a geometric bad-pair fraction $\delta_g$ give
\begin{align}
 \widehat\E\left[
   1_{\rm wrong}e(|F(U)|,\ell)
                      1_{|F(U)|\le R/4}\right]
    &\le C+CR^p\delta_g,\label{eq:weighted-empirical-cost}\\
 \widehat\E[1_{\rm wrong}1_{|F(U)|>R/4}]
    &\le CR^{-p}+\delta_g.\label{eq:weighted-empirical-huge}
\end{align}
Here $\widehat\E$ includes all auxiliary averages. The contact-quantile
exception has $|F|=1$, so its empirical contribution is at most one;
no empirical smallness is inferred from its population quantile alone.

These are bounded head-only losses after the geometric inequalities
have been established. For a fixed contact random seed, the contact
and its distance are fixed within each face sign/order cell.
For each fixed $\ell$, the score $e(b,\ell)$ is strictly increasing
in $b>0$. A subgraph threshold therefore reduces to a comparison of
$b$ with a fixed number depending on $\ell$ and the threshold. After
multiplying by the positive normalization scale, this is an affine
comparison in the head coefficients. Thus the same
arrangement and VC-entropy argument gives dimension at most
$CD_r(m+\log(D_r+2))$. Auxiliary averaging preserves the bound.
Choose $R=\rho^{-t_0}$ after obtaining $\delta_g\le\rho^g$, with
$pt_0<\min(g,1/800)/2$. Uniform transfer of
\eqref{eq:weighted-empirical-cost}--\eqref{eq:weighted-empirical-huge}
then gives a bounded population contact cost and a positive-power
bound on the mass of huge wrong amplitudes. Only these bounded
head-only losses are transferred; the geometric bad relation is not.

\subsection{Passing to an accurate face}
The contact quantile $\zeta_0$ is a fixed sufficiently small constant.
Projection to the prescribed sign half-lines changes the normalized
head by at most $C\sqrt{\zeta_0}+o(1)$ in $L^2$; the fixed term need
only lie below the tolerance in
Lemma~\ref{lem:stable-contact-calibration}. The random-restriction
leakage still tends to zero. A full-head error at most $3/8$ implies
that a positive fraction of faces have error at most $7/16$. On such
a face both sign classes have positive constant mass. Markov selection
for the bounded contact cost and the leakage therefore supplies an
eligible accurate face with cost bounded by a fixed constant.

The huge wrong-amplitude set has positive-power mass by
\eqref{eq:weighted-empirical-huge} and bounded statistical transfer.
If the restricted polynomial has degree $u\le a_1H$, with $a_1$
sufficiently small, hypercontractivity gives precisely the exceptional-set
absorption hypothesis of Lemma~\ref{lem:stable-contact-calibration}.
Its contact-distance hypothesis is
\eqref{eq:contact-distance-moments}. The full contact score gives a stronger moment bound than its
weighted-$p$ simplification. Indeed, on retained pairs,
\[
 |F(U)|\le C_\beta\big\{
 e(|F(U)|,\ell)^{1/2}(\ell+1)^{\beta/2}
       +e(|F(U)|,\ell)^{1/p}\big\}.
\]
Since $2-p=\beta$, H\"older's inequality gives
\[
 \|e^{1/2}(\ell+1)^{\beta/2}\|_p
 \le (\E e)^{1/2}\|\ell+1\|_p^{\beta/2}
 \le C m^{\beta/4}.
\]
The second term has $L^p$ norm $(\E e)^{1/p}$. Thus the wrong-side
$L^p$ norm is $O(m^{\beta/4})$, with the same exceptional-set absorption
as above. No exponent below $p$ or endpoint logarithm is needed.
Theorem~\ref{thm:endpoint-contact-calibration} consequently gives
\begin{equation}\label{eq:balanced-endpoint-degree}
 u\ge c m^{\theta_*},\qquad
 \theta_* =\frac{1-\beta/4}{s_p},\quad p=2-\beta.
\end{equation}
The analytic theorem is separate from the geometric and statistical
reductions above.

Choosing $u=a_1H$ and inverting
\eqref{eq:balanced-endpoint-degree} gives a face size of order
$H^{1/\theta_*}$. The resulting random-face sample exponent is
\begin{equation}\label{eq:balanced-sample-scale}
 kH^{\nu_*},\qquad \nu_*=2-\frac1{\theta_*}.
\end{equation}
The complementary full-face scale is $k^{\theta_*}H$.
For $\beta=1/2$, $\theta_*=0.719367150047\ldots$ and
$\nu_*=0.609889261784\ldots$.

These calibrated exponents use effective profiles at the canonical
coupled contacts. Their actual-posterior application uses the coupled
middle-band, angular, and endpoint estimates proved above; bounded
contact marginals alone would not suffice.

% Standalone analytic module. Requires the main paper's standard theorem environments.
% Bibliography additions are in analytic-calibration-bibliography.tex.
\providecommand{\norm}[2]{\lVert #1\rVert_{#2}}
\section{Analytic calibration of contact-weighted errors}
Let $\mu_m$ be uniform probability measure on $\{-1,1\}^m$. Write
\[
 f(x)=\prod_{i=1}^m x_i,\qquad
 \partial_i P(x)=\frac{P(x)-P(x^{\oplus i})}{2},\qquad
 \Delta=\sum_{i=1}^m\partial_i.
\]
Thus $\Delta\chi_S=|S|\chi_S$. All unqualified norms and expectations use $\mu_m$. For $1<t<2$, put
\[
 s_t=1+\frac{2}{\pi}\arcsin\!\left(\frac2t-1\right).
\]
We use the dimension-free scalar Bernstein inequality
\begin{equation}\label{ac:eq:bernstein}
 \norm{\Delta P}{t}\le C_t u^{s_t}\norm{P}{t}
 \qquad (\deg P\le u).
\end{equation}
This is part of the established theory of polynomial inequalities on the cube \cite{EI}. Section~\ref{ac:sec:log} supplies the needed version from the imaginary-power estimate of Carbonaro and Dragi\v cevi\'c \cite{CD}.

Fix $0<\beta\le1/2$ and put $p=2-\beta$. A contact law is a joint distribution $\pi$ of $(U,V)$ in the cube whose first marginal is exactly $\mu_m$. Set $L=\dist(U,V)$. For a margin-separated function $F$, meaning $|F|\ge1$, define
\[
 E_F=\{x:f(x)F(x)<0\}.
\]
Allow an arbitrary exceptional set $B$ of \emph{pairs} $(U,V)$, and define
\begin{equation}\label{ac:eq:cost}
 \mathcal C_{\pi,B}(F)
 =\E_\pi\!\left[\mathbf 1_{\{U\in E_F,\ (U,V)\notin B\}}
 \min\left\{\frac{|F(U)|^2}{(L+1)^\beta},\ |F(U)|^p\right\}\right].
\end{equation}
In particular, an exceptional set of source points is permitted. No independence between amplitudes, labels, distances, or exceptions is assumed.

\begin{lemma}[Stable contact-cost calibration]\label{lem:stable-contact-calibration}
Fix $c_0>0$, $K<\infty$, and $C_*<\infty$. Suppose that:
\begin{enumerate}
\item $|F|\ge1$ and $\E[f\sign F]\ge c_0$;
\item $Q$ has degree at most $u\ge1$ and $\norm{F-Q}{p}\le e\le c_0/2$;
\item $\pi(B)\le\delta$ and
\begin{equation}\label{ac:eq:absorb}
 \delta^{1/p-1/2}(p-1)^{-u/2}\le\frac1{16};
\end{equation}
\item $\mathcal C_{\pi,B}(F)\le K$ and
\begin{equation}\label{ac:eq:distance}
 \norm{L+1}{L^p(\pi)}\le C_*\sqrt m.
\end{equation}
\end{enumerate}
Then
\begin{equation}\label{ac:eq:theta}
 u\ge c\,m^{\theta_\beta},\qquad
 \theta_\beta=\frac{1-\beta/4}{s_p},
\end{equation}
where $c>0$ depends only on the displayed fixed parameters.
\end{lemma}

The exceptional set may depend arbitrarily on $F,Q$ and the contact law. Its small probability is not sufficient on its own; condition~\eqref{ac:eq:absorb} controls its amplitude contribution. For $\delta=\rho^b$, it suffices to have $u\le a\log(1/\rho)$ with $a$ sufficiently small.

\begin{theorem}[Explicit contact exponent]\label{thm:endpoint-contact-calibration}
The exponent in Lemma~\ref{lem:stable-contact-calibration} is
\begin{equation}\label{ac:eq:endpoint}
 \theta_\beta=\frac{1-\beta/4}
 {1+(2/\pi)\arcsin(\beta/(2-\beta))}.
\end{equation}
For $\beta=1/2$, it equals $0.7193671500\ldots$. There is no endpoint logarithmic loss: the single exponent $p=2-\beta$ is used throughout.
\end{theorem}

\subsection{Proof of the calibration theorem}
We first record two simple consequences of parity.

\begin{lemma}[Parity and positivity]\label{ac:lem:parity}
If $|F|\ge1$, $g=\sign F$, and $\E[fg]\ge c_0$, then
\begin{equation}\label{ac:eq:influence}
 \E[g\Delta F]\ge I(g)\ge mc_0,
 \qquad I(g)=\sum_i\Pp\{g(X)\ne g(X^{\oplus i})\}.
\end{equation}
For every real $Q$ and $1\le t<\infty$,
\begin{equation}\label{ac:eq:positive}
 \norm{Q}{t}\le 2\norm{(fQ)_-}{t}+\frac2m\norm{\Delta Q}{t}.
\end{equation}
\end{lemma}
\begin{proof}
On an edge where $g$ changes sign, margin separation gives
$(g(x)-g(x^{\oplus i}))(F(x)-F(x^{\oplus i}))\ge4$.
Edge symmetrization proves the first inequality in~\eqref{ac:eq:influence}.
Also $\E[fg]=\E[f\partial_i g]$, so
$|\E[fg]|\le\E|\partial_i g|=\operatorname{Inf}_i(g)$ for every $i$.
Summing proves its second inequality.

Let $A$ be normalized cube adjacency and $q=fQ$. Since parity reverses on every edge,
$(I+A)q=2f\Delta Q/m$. Write its right side as $h$. Positivity of $A$ gives
$q_+\le A q_-+|h|$. Contractivity and the triangle inequality give
$\norm{q}{t}\le2\norm{q_-}{t}+\norm{h}{t}$.
\end{proof}

\begin{proof}[Proof of Lemma~\ref{lem:stable-contact-calibration}]
Set $X=\norm Qp$, and let
\[
 b=|F(U)|\mathbf1_{\{U\in E_F,\ (U,V)\notin B\}},\qquad
 z=\min\{b^2/(L+1)^\beta,b^p\},\qquad W=\norm b{L^p(\pi)}.
\]
Pointwise,
$b\le z^{1/2}(L+1)^{\beta/2}+z^{1/p}$.
Since $2-p=\beta$, H\"older's inequality with exponents $2/p$ and $2/\beta$ yields
\begin{align}
 \norm{z^{1/2}(L+1)^{\beta/2}}{L^p(\pi)}
 &\le(\E_\pi z)^{1/2}\norm{L+1}{L^p(\pi)}^{\beta/2},\label{ac:eq:holder}\\
 W&\le K^{1/2}C_*^{\beta/2}m^{\beta/4}+K^{1/p}
 \le C m^{\beta/4}.\label{ac:eq:W}
\end{align}
Keeping the quadratic part of the contact cost is essential for this bound.

The exact uniform first marginal and H\"older's inequality give
\[
 \norm{Q(U)\mathbf1_B}{L^p(\pi)}
 \le\delta^{1/p-1/2}\norm Q2
 \le\delta^{1/p-1/2}(p-1)^{-u/2}X\le X/16,
\]
where the second inequality is the degree-$u$ hypercontractive comparison.
It follows that
$\norm{F\mathbf1_{E_F}}p\le W+X/16+e$.
The negative-part map is $1$-Lipschitz, so Lemma~\ref{ac:lem:parity} implies
\begin{equation}\label{ac:eq:boot}
 X\le2W+X/8+4e+\frac2m\norm{\Delta Q}p.
\end{equation}
If $2C_pu^{s_p}/m\ge1/4$, then $u\ge c m^{1/s_p}$, which is stronger than the conclusion. Otherwise~\eqref{ac:eq:bernstein} absorbs the last term in~\eqref{ac:eq:boot}, giving $X\le C(W+e)$.

Each $\partial_i$ is an $L^1$ contraction, so
$\norm{\Delta(F-Q)}1\le me$. By~\eqref{ac:eq:influence},
\[
 \norm{\Delta Q}p\ge\E[g\Delta Q]\ge m(c_0-e)\ge mc_0/2.
\]
Combining this with~\eqref{ac:eq:bernstein},~\eqref{ac:eq:W}, and the bound on $X$ proves~\eqref{ac:eq:theta}. Formula~\eqref{ac:eq:endpoint} is the substitution $p=2-\beta$.
\end{proof}

\section{A logarithmic Bernstein refinement}\label{ac:sec:log}
The next estimate makes the norm dependence in scalar Bernstein explicit. Its proof follows the fixed-vector interpolation method of Domelevo et al.\ \cite{Domelevo}, with strip harmonic measure in place of Gaussian damping. We make no priority claim for the interpolation refinement.

\begin{theorem}\label{ac:thm:log}
For $1<t<2$, set $\sigma_t=(2/\pi)\arcsin(2/t-1)$. Every nonzero real cube polynomial of degree at most $u\ge1$ satisfies
\begin{equation}\label{ac:eq:logbern}
 \norm{\Delta P}t\le C_t u
 \left[1+\log\frac{\norm P2}{\norm Pt}\right]^{\sigma_t}\norm Pt.
\end{equation}
The constants are bounded when $t$ ranges over a compact subinterval of $(1,2)$.
\end{theorem}
\begin{proof}
Let $0<\varepsilon\le(t-1)/t$ and define
\[
 q=\frac{t(1-\varepsilon)}{1-t\varepsilon/2},\qquad
 \frac1t=\frac\varepsilon2+\frac{1-\varepsilon}{q}.
\]
We first prove
\begin{equation}\label{ac:eq:mixed}
 \norm{\Delta P}t\le C_tu\varepsilon^{-\sigma_t}
 \norm P2^\varepsilon\norm Pq^{1-\varepsilon}.
\end{equation}
For $L=\Delta+\gamma I$, $\gamma>0$, the imaginary-power estimate of \cite[Proposition 11]{CD} is
\[
 \norm{L^{iv}}{q\to q}\le C(q^*-1)(1+|v|)^{1/2}
 e^{\phi_q|v|},\qquad \phi_q=\arcsin(2/q-1).
\]
Its constants are independent of the dimension and of $\gamma$. Fix a real $g$ with $\norm g{t'}=1$, put $N=1/\varepsilon$, and form the holomorphic bilinear pairing
\[
 G(z)=\E\left[(L^{N(1-z)}P)\,g_z\right],\qquad
 g_z=\sign(g)|g|^{t'((1-z)/2+z/q')}.
\]
The value is set to zero where $g=0$. At $\theta=1-\varepsilon$, it equals $\E[(LP)g]$. On the left boundary,
$|G(iv)|\le(u+\gamma)^N\norm P2$; on the right,
\[
 |G(1+iv)|\le C(q^*-1)(1+N|v|)^{1/2}
 e^{\phi_qN|v|}\norm Pq.
\]
On the finite cube $G$ is bounded on the closed strip, so the Poisson inequality applies to $\log|G|$.

The right-boundary harmonic-measure density at $\theta$ is
\[
 w_\varepsilon(v)=\frac{\sin(\pi\varepsilon)}{2(\cosh(\pi v)-\cos(\pi\varepsilon))},
 \qquad \int w_\varepsilon=1-\varepsilon.
\]
The elementary bounds needed here are
\begin{align*}
 \int |v|w_\varepsilon(v)\,dv
 &\le\frac{2\varepsilon}{\pi}\bigl(\log(1/\varepsilon)+1\bigr),\\
 \int\log(1+|v|/\varepsilon)w_\varepsilon(v)\,dv&\le C.
\end{align*}
For the first, expand
\[\frac{\sin a}{\cosh x-\cos a}=2\sum_{n\ge1}\sin(na)e^{-nx}\]
and integrate, obtaining
\[(2/\pi^2)\int_0^{\pi\varepsilon}-\log(2\sin(s/2))\,ds.\]
For the second, use $w_\varepsilon(v)\le\pi\varepsilon/[4(v^2+\varepsilon^2)]$ and rescale.

The left boundary has mass $\varepsilon$. Integrating the boundary logarithms consequently gives
\[
 |G(\theta)|\le C_t(u+\gamma)
 \norm P2^\varepsilon\norm Pq^{1-\varepsilon}
 \varepsilon^{-2\phi_q/\pi}.
\]
Here $q\ge1+(t-1)/2$, and $|\phi_q-\phi_t|\le C_t\varepsilon$, so the last power is at most $C_t\varepsilon^{-\sigma_t}$. Taking the dual supremum and then $\gamma\downarrow0$ proves~\eqref{ac:eq:mixed}.

Since $q<t$, $\norm Pq\le\norm Pt$. Take
$\varepsilon=\min\{(t-1)/t,[1+\log(\norm P2/\norm Pt)]^{-1}\}$ to obtain~\eqref{ac:eq:logbern}. All bounds used above are uniform on compact exponent intervals.
\end{proof}

Hypercontractivity gives $\log(\norm P2/\norm Pt)\le(u/2)\log(1/(t-1))$. Substitution recovers~\eqref{ac:eq:bernstein}, since $s_t=1+\sigma_t$. Theorem~\ref{ac:thm:log} therefore does not prove the general linear-degree Bernstein conjecture; it gives linear degree when the norm ratio is bounded.


\section{Random restriction to an accurate face}\label{sec:face-restriction}
For a random $m_f$-face and its degree-$u$ Fourier truncation,
\begin{equation}\label{eq:face-restriction}
 \E_{\rm face}\|P-P_{\rm face,\le u}\|_2^2
 \le\|P\|_2^2
       \left(\frac{e(r/k)m_f}{u+1}\right)^{u+1}.
\end{equation}
To see this, expand in Walsh characters. A support of size at most $r$
retains more than $u$ active coordinates with at most the displayed
hypergeometric upper-tail probability. Random fixed signs cancel all
cross terms. Averaging and summing the squared coefficients proves the
inequality.

If the global head error is at most $3/8+o(1)$, a positive fraction of
faces have error at most $7/16$ and hence correlation at least $1/8$.
Each sign class on such a face has mass at least $1/16$.
The global head leakage and Markov discard only a vanishing fraction
for the bad-margin eligibility condition. The population versions of
\eqref{eq:weighted-empirical-cost}--\eqref{eq:weighted-empirical-huge},
together with~\eqref{eq:face-restriction}, with fixed
Markov constants, select one accurate eligible face with controlled
wrong-side moment, huge-set mass, and restriction leakage.
Normalize its truncation with the same $s,t$:
$P_0=(P_{\rm face,\le u}-t)/s$. The contact projection bound gives
\[
 \|F-P_0\|_p\le C\sqrt{\zeta_0}
                +2\|P-P_{\rm face,\le u}\|_2/a.
\]
Take $u=a_1H$ with $a_1$ sufficiently small relative to $b_0$.
Then the huge-set mass $C\rho^{b_0}$ satisfies the absorption condition
in Theorem~\ref{thm:endpoint-contact-calibration}. Its conclusion forces
$u\ge c m_f^\theta$.

\section{Choice of parameters and proof of the main theorem}
\label{sec:assembly}
\begin{proof}[Proof of Theorem~\ref{thm:main}]
Fix small confidence parameters for the comparator, sampling, Gram,
transport, and bounded-loss events so their failure probabilities sum
to at most $1/4$. All conclusions are simultaneous over the minimizing
fields. Suppose one such field has test error less than $3/8$.
Proposition~\ref{prop:head} transfers this error and the global margin
to its low-degree head up to vanishing population errors.

We record one order of parameter choices. Put
$\zeta_{\rm stat}=1/400$, $e=\zeta_{\rm stat}/100$,
$\epsilon=\rho^e$, and use the head cap
$a\rho^{-\beta/2-10e}$. Take recurrent window mass
$\rho^{\beta/2+30e}$, a fixed sufficiently small contact quantile mass $\zeta_0$, and anchor horizon
$S=\rho^{-\Theta}$ with $\Theta=(1-\beta)/2$.
The transport estimates leave a geometric bad fraction at most
$\rho^g$, with $g=5e$ permissible. Choose
\[
 0<t_{\rm rad}<\min\{g/(4p),\zeta_{\rm stat}/(4p)\},
 \qquad R=\rho^{-t_{\rm rad}}.
\]
Then both the geometric bad contribution $R^p\rho^g$ and the bounded
statistical error $R^p\rho^{\zeta_{\rm stat}/2}$ tend to zero.
All fixed powers of $H$ are smaller than $R$.

First consider $k$ at least a sufficiently large multiple of
$H^{1/\theta}$. Choose $u=\lfloor a_1H\rfloor$ with $a_1>0$
sufficiently small for bad-moment absorption. Choose a large fixed
$C_2$ and let $m_f=\lceil C_2u^{1/\theta}\rceil$.
Finally choose the degree ratio so small that
\[
 r/k\le c_2u/m_f.
\]
The head norm bound and~\eqref{eq:face-restriction} make restriction
leakage a positive power of $\rho$ after decreasing $c_2$.
Theorem~\ref{thm:endpoint-contact-calibration} forces
$u\ge c m_f^\theta$, contrary to
the choice of $C_2$. Since $H\asymp\log(ed/k)$,
\[
 rH\asymp kH^{2-1/\theta}.
\]
The constants can be selected within fixed factors of the cutoff in
\eqref{eq:cutoff-choice}, with harmless odd-integer rounding.

For $k$ below a sufficiently large multiple of $H^{1/\theta}$, use the
entire active cube as the face: $m_f=k$ and $u=r$.
There is no restriction leakage. Choose an odd
$r\asymp c_0 k^\theta$, with $c_0$ small enough that $r\le a_1H$.
The same bad-moment absorption and calibration yield
$r\ge c k^\theta$, a contradiction. Here
$\log n_r\asymp k^\theta H$. The two branches overlap after adjusting
fixed constants and agree in scale at $k\asymp H^{1/\theta}$.
In both branches $k=O(rH)$, $m_f$ is polynomial in $H$, and
$\log n_r\asymp rH=o(d)$, as required by the preceding estimates.
Thus every compatible minimizing sign has test error at least $3/8$.

It remains to pass to the actual posterior. Proposition~\ref{prop:rate}
puts asymptotic posterior mass one on the minimum-energy compatible
signs at every fixed $0<\gamma<1$. For any mixture among these signs,
let $q(x)$ be its correct-label probability and
$R_{\rm draw}=\E_x[1-q(x)]$. Then $R_{\rm draw}\ge3/8$.
If $F=\Pp_x\{q(x)\le2/3\}$,
\[
 R_{\rm draw}\le(1-F)/3+F,
 \qquad F\ge(3R_{\rm draw}-1)/2\ge1/16.
\]
The finite-width draw error is at least $3/8-o(1)$; applying the same
inequality before taking the lower limit proves the theorem without
assuming convergence of the relative weights among tied minimizers.
\end{proof}



\section{Scope}
The result concerns the specified globally conditioned posterior, rather
than a posterior mean or a kernel ridge predictor. The confidence
criterion is applied to the actual label probability. The lower bound
holds for every fixed $0<\gamma<1$ because the same finite-dimensional
kernel rate governs this entire interval when width is taken first.
No assertion about $\gamma=0$ is made here.

For fixed polynomial sparsity, the natural remaining target is
$\exp\{c k\log(ed/k)\}$. The exponent in this paper is smaller. The
improvement comes from using the global margin during the proof, rather
than using only the low-degree head's norm. A matching natural-scale
lower bound would require a stronger local parity calibration or a
different variational argument.

\begin{thebibliography}{99}
\bibitem{flajolet2009}
P. Flajolet and R. Sedgewick. \emph{Analytic Combinatorics}.
Cambridge University Press, 2009. Chapter VI.
\bibitem{Lee2023}
J. R. Lee. Spectral hypergraph sparsification via chaining.
\href{https://arxiv.org/abs/2209.04539}{arXiv:2209.04539}, Section 2.4.
\bibitem{tropp2008}
J. A. Tropp. Norms of random submatrices and sparse approximation.
\emph{Comptes Rendus Math\'ematique} 346:1271--1274, 2008.
\bibitem{tropp2012}
J. A. Tropp. User-friendly tail bounds for sums of random matrices.
\emph{Foundations of Computational Mathematics} 12:389--434, 2012.
\bibitem{EI}
A. Eskenazis and P. Ivanisvili.
Polynomial inequalities on the Hamming cube.
\href{https://arxiv.org/abs/1902.02406}{arXiv:1902.02406}.
\bibitem{HKS}
A. W. Harrow, A. Kolla, and L. J. Schulman.
Dimension-free $L^2$ maximal inequality for spherical means in the
hypercube. \emph{Theory of Computing} 10:55--75, 2014.
\bibitem{CD} A. Carbonaro and O. Dragi\v cevi\'c.
\emph{Functional calculus for generators of symmetric contraction semigroups}.
Duke Math. J. 166 (2017), 937--974.
\href{https://arxiv.org/abs/1308.1338}{arXiv:1308.1338}; imaginary-power estimate in Proposition 11 of version 2.
\bibitem{Domelevo} K. Domelevo, P. Durcik, V. Fragkiadaki, O. Klein, D. Oliveira e Silva, L. Slav\'ikov\'a, and B. Wr\'obel.
\emph{Dimension-free estimates for low degree functions on the Hamming cube}.
Studia Math. 280 (2025), 161--173.
\href{https://arxiv.org/abs/2401.07699}{arXiv:2401.07699}.

\end{thebibliography}
\end{document}
